\section{Technical examples and source-specific certificates}
This appendix collects finite certificates, source-specific endpoint checks, numerical witnesses, and applications whose details are useful for verification but would interrupt the main narrative.
\begin{counterexample}
\label{res:baskakov-alpha1-r8-not-cm}
The Baskakov \(\alpha=1\), \(r=8\) function \(f^{[8]}_1(x)=1/((1+x)^8-x^8)\) is not completely monotone on \((0,\infty)\).
\end{counterexample}
\begin{proof}
Set
\[
D_m(x)=(1+x)^{2m}-x^{2m}.
\]
The roots are obtained from
\[
\frac{1+x}{x}=\omega_k,\qquad
\omega_k=e^{\pi i k/m},\qquad k=1,\ldots,2m-1.
\]
Thus
\[
\lambda_{k,m}
=\frac{1}{\omega_k-1}
=-\frac12-\frac{i}{2}\cot\frac{\pi k}{2m}.
\]
At this root,
\[
D_m'(\lambda_{k,m})
=2m\lambda_{k,m}^{2m-1}(\omega_k^{-1}-1),
\]
and direct simplification gives the real residue
\[
R_{k,m}
=\frac{1}{D_m'(\lambda_{k,m})}
=\frac{(-1)^{m+k}}{2m}
\left(2\sin\frac{\pi k}{2m}\right)^{2m-2}.
\]
Consequently the inverse-Laplace density is
\[
\rho_m(t)
=\sum_{k=1}^{2m-1}R_{k,m}e^{\lambda_{k,m}t}.
\]
Pairing conjugate roots gives
\[
\rho_m(t)
=\frac{e^{-t/2}}{2m}
\left[
2^{2m-2}
+2\sum_{k=1}^{m-1}
(-1)^{m+k}
\left(2\sin\frac{\pi k}{2m}\right)^{2m-2}
\cos\!\left(
\frac{t}{2}\cot\frac{\pi k}{2m}
\right)
\right].
\]

This proves the finite trigonometric density formula used by the diagnostic
candidate.

For \(m=2\), this gives
\[
\rho_2(t)=e^{-t/2}\left(1-\cos\frac t2\right)
=2e^{-t/2}\sin^2\frac t4\ge0,
\]
recovering the previously admitted \(r=4,\alpha=1\) seed.

For \(m=3\),
\[
\rho_3(t)
=\frac{e^{-t/2}}6
\left[
16+2\cos\left(\frac{\sqrt3\,t}{2}\right)
-18\cos\left(\frac{t}{2\sqrt3}\right)
\right].
\]
Set \(u=t/(2\sqrt3)\).  Since
\(\cos(3u)=4\cos^3u-3\cos u\),
\[
\rho_3(t)
=\frac{4}{3}e^{-t/2}(1-\cos u)^2(2+\cos u)\ge0.
\]
Thus \(r=6,\alpha=1\) is completely monotone by an explicit positive density.

For \(m=4\), write \(h_4(t)=e^{t/2}\rho_4(t)\).  The density formula gives
\[
h_4(t)
=8+2\cos\frac t2
-\frac{(2-\sqrt2)^3}{4}
\cos\left(\frac{1+\sqrt2}{2}t\right)
-\frac{(2+\sqrt2)^3}{4}
\cos\left(\frac{\sqrt2-1}{2}t\right).
\]
At \(t=10\pi\),
\[
\cos(t/2)=\cos(5\pi)=-1,
\]
and both other cosine terms equal \(-\cos(5\pi\sqrt2)\).  Since
\[
(2-\sqrt2)^3+(2+\sqrt2)^3=40,
\]
we obtain
\[
h_4(10\pi)=6+10\cos(5\pi\sqrt2).
\]
Now
\[
0<5\sqrt2-7<\frac14
\]
because \(7/5<\sqrt2<29/20\).  Hence, with
\(\delta=5\sqrt2-7\),
\[
\cos(5\pi\sqrt2)=\cos(7\pi+\pi\delta)=-\cos(\pi\delta)
<-\cos\frac\pi4=-\frac{\sqrt2}{2}<-\frac35.
\]
Therefore
\[
h_4(10\pi)<6+10\left(-\frac35\right)=0,
\]
and
\[
\rho_4(10\pi)=e^{-5\pi}h_4(10\pi)<0.
\]
Since the inverse-Laplace transform is unique for measures on
\([0,\infty)\), \(f^{[8]}_1(x)=1/((1+x)^8-x^8)\) cannot be completely monotone.
\end{proof}

\begin{counterexample}
\label{res:qi-hlambda-x4-source-range-not-cm}
For \(h_\lambda(x)=\Psi(x)-(x^2+\lambda x+12)/(12x^4(x+1)^2)\), \(x^4h_\lambda(x)\) is not completely monotone for every \(\lambda\le0\), and \(x^4(-h_\mu(x))\) is not completely monotone for every \(\mu\ge4\).
\end{counterexample}
\begin{proof}
A completely monotone function must be nonincreasing, so its first derivative cannot be positive at any point.

Use the recurrence expansions at \(x=0^+\):
\[
\psi'(x)=\frac1{x^2}+\zeta(2)-2\zeta(3)x+3\zeta(4)x^2+O(x^3),
\]
and
\[
\psi''(x)=-\frac2{x^3}-2\zeta(3)+6\zeta(4)x+O(x^2).
\]
Therefore
\[
\Psi(x)
=\frac1{x^4}-\frac2{x^3}+\frac{\pi^2}{3x^2}
-\frac{4\zeta(3)}x+O(1).
\]
Also
\[
\frac{x^2+\lambda x+12}{12x^4(1+x)^2}
=\frac1{x^4}+\left(\frac{\lambda}{12}-2\right)\frac1{x^3}
+\left(\frac{37}{12}-\frac{\lambda}{6}\right)\frac1{x^2}
+O(x^{-1}).
\]
Subtracting gives
\[
h_\lambda(x)
=-\frac{\lambda}{12x^3}
+\left(\frac{\lambda}{6}+\frac{\pi^2}{3}-\frac{37}{12}\right)\frac1{x^2}
+O(x^{-1}).
\]

If \(\lambda<0\), then
\[
x^4h_\lambda(x)=-\frac{\lambda}{12}x+O(x^2),
\]
so
\[
\frac{d}{dx}\bigl[x^4h_\lambda(x)\bigr]
=-\frac{\lambda}{12}+O(x)>0
\]
for all sufficiently small \(x>0\). Hence \(x^4h_\lambda\) is not completely monotone.

If \(\lambda=0\), then
\[
x^4h_0(x)
=\left(\frac{\pi^2}{3}-\frac{37}{12}\right)x^2+O(x^3).
\]
Since \(4\pi^2>37\), the leading coefficient is positive, and
\[
\frac{d}{dx}\bigl[x^4h_0(x)\bigr]
=2\left(\frac{\pi^2}{3}-\frac{37}{12}\right)x+O(x^2)>0
\]
for all sufficiently small \(x>0\). Hence \(x^4h_0\) is not completely monotone.

Similarly, for \(\mu\ge4\),
\[
x^4(-h_\mu(x))=\frac{\mu}{12}x+O(x^2),
\]
so
\[
\frac{d}{dx}\bigl[x^4(-h_\mu(x))\bigr]
=\frac{\mu}{12}+O(x)>0
\]
for all sufficiently small \(x>0\). Hence \(x^4(-h_\mu)\) is not completely monotone for every \(\mu\ge4\).

Thus the degree-four transforms required by Qi's conjecture are not completely monotone on the conjectured source ranges. The conjecture in Remark 7.4 is false.
\end{proof}

\begin{proposition}
\label{res:keady-inverse-third-derivative-sign-certificate}
For \(f(x)=x^{-1}+100e^{-x}\) and \(g=f^{-1}\), the inverse-derivative certificate gives \(g'''(f(1/8))>0\).
\end{proposition}
\begin{proof}
There exists a completely monotone decreasing bijection
\[
f:(0,\infty)\to(0,\infty)
\]
whose inverse is not completely monotone. One explicit example is
\[
f(x)=\frac1x+100e^{-x}.
\]

For every \(n\ge0\),
\[
(-1)^n f^{(n)}(x)=\frac{n!}{x^{n+1}}+100e^{-x}>0.
\]
Thus \(f\) is strictly completely monotone. Also
\[
f'(x)=-x^{-2}-100e^{-x}<0,
\]
so \(f\) is strictly decreasing. Finally,
\[
\lim_{x\downarrow0}f(x)=\infty,
\qquad
\lim_{x\to\infty}f(x)=0.
\]
Therefore \(f\) is a decreasing bijection from \((0,\infty)\) onto \((0,\infty)\).

Let \(g=f^{-1}\) and write \(y=f(x)\). Inverse differentiation gives
\[
g'''(f(x))=\frac{3(f''(x))^2-f'''(x)f'(x)}{(f'(x))^5}.
\]
For \(f_a(x)=x^{-1}+ae^{-x}\), define
\[
N_a(x)=3(f_a''(x))^2-f_a'''(x)f_a'(x).
\]
A direct calculation gives
\[
N_a(x)=\frac6{x^6}+ae^{-x}\left(-\frac6{x^4}+\frac{12}{x^3}-\frac1{x^2}\right)+2a^2e^{-2x}.
\]
At \(a=100\) and \(x_0=1/8\),
\[
N_{100}(1/8)=1{,}572{,}864-1{,}849{,}600e^{-1/8}+20{,}000e^{-1/4}.
\]
Using \(e^{-1/8}>7/8\) and \(e^{-1/4}<1\),
\[
N_{100}(1/8)<1{,}572{,}864-1{,}849{,}600\cdot\frac78+20{,}000=-25{,}536<0.
\]
Since \(f'(1/8)<0\), the denominator \((f'(1/8))^5\) is negative. Hence
\[
g'''(f(1/8))=\frac{N_{100}(1/8)}{(f'(1/8))^5}>0.
\]
A completely monotone function must satisfy \((-1)^3g'''\ge0\), i.e. \(g'''\le0\). Therefore \(g=f^{-1}\) is not completely monotone.
\end{proof}

\begin{proposition}
\label{res:yin-pk-digamma-near-zero-asymptotic}
For fixed \(p\in\mathbb N\) and \(k>0\), the Yin bracket satisfies \(B_{p,k}(x)=x^{-1}+O(|\log x|)\) as \(x\to0^+\). Consequently \(x^\alpha B_{p,k}(x)\sim x^{\alpha-1}\).
\end{proposition}
\begin{proof}
\emph{Setup.}
Fix \(p\in\mathbb N\) and \(k>0\). Define
\[
B_{p,k}(x)=
\frac1k\log\frac{pkx}{x+k(p+1)}-\psi_{p,k}(x),
\qquad
\delta_{p,k,\alpha}(x)=x^\alpha B_{p,k}(x).
\]

The source gives the finite formula
\[
\psi_{p,k}(x)=\frac1k\log(pk)-\sum_{n=0}^p\frac1{x+nk}.
\]
Therefore
\[
\begin{aligned}
B_{p,k}(x)
&=
\frac1k\log\frac{pkx}{x+k(p+1)}
-\frac1k\log(pk)
+\sum_{n=0}^p\frac1{x+nk}  \\
&=
\frac1k\log\frac{x}{x+k(p+1)}
+\sum_{n=0}^p\frac1{x+nk}.
\end{aligned}
\]

This normal form also gives positivity directly. Put \(t=x/k\). Then
\[
B_{p,k}(x)
=\frac1k\left[
\sum_{n=0}^p\frac1{t+n}
+\log\frac{t}{t+p+1}
\right].
\]
Since
\[
\log\frac{t+p+1}{t}=\int_0^{p+1}\frac{du}{t+u},
\]
and \(u\mapsto(t+u)^{-1}\) is strictly decreasing, the left Riemann sum dominates:
\[
\sum_{n=0}^p\frac1{t+n}
>
\int_0^{p+1}\frac{du}{t+u}.
\]
Hence \(B_{p,k}(x)>0\) for \(x>0\).

This proves the claim.

From the finite normal form,
\[
B_{p,k}(x)
=\frac1x
+\frac1k\log\frac{x}{x+k(p+1)}
+\sum_{n=1}^p\frac1{x+nk}.
\]

As \(x\to0^+\), the finite sum over \(n\ge1\) is \(O(1)\), and
\[
\frac1k\log\frac{x}{x+k(p+1)}=O(|\log x|).
\]
Thus
\[
B_{p,k}(x)=\frac1x+O(|\log x|)
\qquad (x\to0^+).
\]
Equivalently,
\[
xB_{p,k}(x)\to1.
\]
Therefore
\[
\delta_{p,k,\alpha}(x)
=x^\alpha B_{p,k}(x)
=x^{\alpha-1}(xB_{p,k}(x))
\sim x^{\alpha-1}.
\]

This proves the claim.

Assume that \(\delta_{p,k,\alpha}\) is completely monotone on \((0,\infty)\). Then
\[
\delta_{p,k,\alpha}(x)\ge0,\qquad
\delta_{p,k,\alpha}'(x)\le0.
\]
Since \(B_{p,k}(x)>0\), the function \(\delta_{p,k,\alpha}\) is actually positive for every \(x>0\).

Suppose, for contradiction, that \(\alpha>1\). The endpoint asymptotic gives
\[
\lim_{x\to0^+}\delta_{p,k,\alpha}(x)=0.
\]
Fix \(x_0>0\). Positivity gives \(\delta_{p,k,\alpha}(x_0)>0\). Because \(\delta_{p,k,\alpha}(x)\to0\) as \(x\to0^+\), choose \(0<y<x_0\) with
\[
\delta_{p,k,\alpha}(y)<\frac12\delta_{p,k,\alpha}(x_0).
\]
But \(\delta_{p,k,\alpha}'\le0\) means the function is nonincreasing as \(x\) increases, hence for \(0<y<x_0\),
\[
\delta_{p,k,\alpha}(y)\ge\delta_{p,k,\alpha}(x_0),
\]
a contradiction.

Therefore \(\alpha\le1\).

This proves the claim.
\end{proof}

\begin{proposition}
\label{res:ramanujan-turan-n2-upper-window-counterfamily}
For the Ramanujan integral Turan gap, there are parameters \(\alpha\in(0,1/2)\) such that \(H_2(x;\alpha)<0\) for some \(x>0\); specifically, with \(x=e^{-L}\) and \(\alpha_L=1/2-1/(4L)\), this holds for all sufficiently large \(L\).
\end{proposition}
\begin{proof}
Put
\[
A_k(x)=(-1)^k I_R^{(k)}(x)
=\int_0^\infty e^{-xt}\frac{t^{k-1}}{\pi^2+\log^2 t}\,dt
\qquad (k\ge1).
\]
Then
\[
H_2(x;\alpha)=A_2(x)^2-\alpha A_1(x)A_3(x).
\]

We study \(x=e^{-L}\) as \(L\to\infty\). With \(y=xt\),
\[
A_k(e^{-L})
=
e^{kL}\int_0^\infty
\frac{e^{-y}y^{k-1}}
{\pi^2+(L+\log y)^2}\,dy.
\]

For fixed \(k=1,2,3\),
\[
A_k(e^{-L})
=
e^{kL}L^{-2}\Gamma(k)
\left(1-\frac{2\psi(k)}{L}+O(L^{-2})\right),
\]
where \(\psi\) is the digamma function. To justify the expansion, split the \(y\)-integral into the central region \(|\log y|\le L^{2/3}\) and its complement. On the central interval, with \(z=\log y\), expand uniformly:
\[
\frac{1}{\pi^2+(L+\log y)^2}
=
L^{-2}\left(1-\frac{2\log y}{L}+O\left(\frac{1+\log^2 y}{L^2}\right)\right),
\]
and integrate against \(e^{-y}y^{k-1}\). The lower complement \(0<y<e^{-L^{2/3}}\) is \(O(e^{-cL^{2/3}})\) for \(k=1,2,3\) after the change \(y=e^{-u}\), and the upper complement \(y>e^{L^{2/3}}\) is super-exponentially small because of \(e^{-y}\). These errors are \(o(L^{-N})\) for every fixed \(N\), which is more than needed for the displayed expansion.

Therefore
\[
\frac{A_2(e^{-L})^2}{A_1(e^{-L})A_3(e^{-L})}
=
\frac{\Gamma(2)^2}{\Gamma(1)\Gamma(3)}
\left(
1+\frac{2(\psi(1)+\psi(3)-2\psi(2))}{L}+O(L^{-2})
\right).
\]
Using
\[
\psi(1)=-\gamma,\qquad
\psi(2)=1-\gamma,\qquad
\psi(3)=\frac32-\gamma,
\]
we get
\[
\psi(1)+\psi(3)-2\psi(2)=-\frac12,
\]
so
\[
\frac{A_2(e^{-L})^2}{A_1(e^{-L})A_3(e^{-L})}
=
\frac12-\frac{1}{2L}+O(L^{-2}).
\]

Choose
\[
\alpha_L=\frac12-\frac{1}{4L}.
\]
For all sufficiently large \(L\), \(\alpha_L\in(0,1/2)\) and
\[
\frac{A_2(e^{-L})^2}{A_1(e^{-L})A_3(e^{-L})}<\alpha_L.
\]
Consequently
\[
H_2(e^{-L};\alpha_L)
=
A_2(e^{-L})^2-\alpha_L A_1(e^{-L})A_3(e^{-L})
<0.
\]

Complete monotonicity would imply nonnegativity at order zero. Hence \(H_2(\cdot;\alpha_L)\) is not completely monotone for these fixed parameters \(\alpha_L\in(0,1/2)\). The source's universal Turan-window statement is false already in the \(n=2\) upper window.

The new mechanism is a logarithmic-density moment-ratio asymptotic obstruction for quadratic Laplace-moment Turan gaps. It is distinct from the public APP-0014 Stieltjes/complete-Bernstein classification of \(I_R\).
\end{proof}

\begin{proposition}[Gamma-quotient crossing on \((0,1/8)\)]
\label{res:bz-gamma-quotient-n-single-crossing-one-eight}
Let \(N(x)\) be the Bulboaca--Zayed derivative-sign numerator from \(T\)-BZ-gamma-quotient-derivative-normal-form. On \([1,8]\), \(N\) has exactly one zero \(\xi\), with \(N(x)<0\) for \(1\le x<\xi\) and \(N(x)>0\) for \(\xi<x\le8\).
\end{proposition}
\begin{proof}
\emph{Setup.}
Use the notation from the previous Bulboaca--Zayed pass:
\[
G(x)=\log\Gamma(x+1),\qquad
D(x)=\log\frac{x^2+6}{x+6},
\]
and
\[
F(x)=\frac{G(x)}{D(x)}.
\]
For \(x>1\), \(D(x)>0\), and
\[
D'(x)=p(x)=\frac{x^2+12x-6}{(x^2+6)(x+6)}.
\]
Let
\[
N_0(x)=\psi(x+1)D(x)-G(x)D'(x).
\]
The derivative-sign numerator recorded in
the corresponding result is
\[
N(x)=((x^2+6)(x+6))\,N_0(x),
\]
so \(N\) and \(N_0\) have the same sign on \(x>1\).

The source itself proves the decreasing part on \((-1,1)\), and the earlier
the preceding theorem proves \(F'(x)>0\) for
\(x\ge8\).  It remains to prove a single sign change of \(N_0\) on \((1,8]\).

Put
\[
P(x)=x^2+12x-6,\qquad Q(x)=(x^2+6)(x+6),
\]
so \(p=P/Q\), and define
\[
r(x)=\frac{\psi(x+1)}{p(x)}=\psi(x+1)\frac{Q(x)}{P(x)}.
\]
Then
\[
P(x)^2r'(x)=S(x),
\]
where
\[
S(x)=\psi'(x+1)Q(x)P(x)+\psi(x+1)K(x),
\]
and
\[
K(x)=Q'(x)P(x)-Q(x)P'(x)
=x^4+24x^3+48x^2-144x-468.
\]
Direct differentiation gives
\[
S'(x)=P(x)\{Q(x)\psi''(x+1)+2Q'(x)\psi'(x+1)\}+K'(x)\psi(x+1).
\]

For \(y=x+1\ge2\), use the standard elementary bounds
\[
\psi(y)>\log y-\frac1y,\qquad
\psi'(y)>\frac1y+\frac{1}{2y^2},\qquad
\psi''(y)>-\frac1{y^2}-\frac2{y^3}.
\]
The last inequality follows from the series
\[
\psi''(y)=-2\sum_{n=0}^{\infty}(n+y)^{-3}
\]
and an integral upper bound for the sum.

Substituting these bounds into \(S'\) yields
\[
S'(x)>
P(x)\frac{5x^4+30x^3+57x^2+12x-90}{(x+1)^3}
K'(x)\left(\log(x+1)-\frac1{x+1}\right).
\]
On \(x\ge1\), \(P(x)>0\).  The quartic numerator has value \(14\) at \(x=1\)
and positive derivative
\[
20x^3+90x^2+114x+12.
\]
Also
\[
K'(x)=4x^3+72x^2+96x-144
\]
has value \(28\) at \(x=1\) and positive derivative for \(x\ge1\), while
\(\log(x+1)-1/(x+1)>0\) for \(x\ge1\).  Hence
\[
S'(x)>0\qquad(1\le x\le8).
\]

At the endpoints,
\[
S(1)=343\left(\frac{\pi^2}{6}-1\right)-539(1-\gamma)<0,
\]
for example using \(\pi^2<987/100\) and \(\gamma<5773/10000\), which gives
\(-66003/10000<0\).  Also \(S(8)>0\), because \(K(8)=17836>0\),
\(\psi'(9)>0\), and \(\psi(9)=H_8-\gamma>0\).  Therefore \(S\) has exactly
one zero on \([1,8]\).  Since \(P^2>0\), \(r'\) has exactly one zero; \(r\)
decreases first and then increases.

Now define
\[
I(x)=D(x)r(x)-G(x).
\]
Then
\[
N_0(x)=p(x)I(x),\qquad I'(x)=D(x)r'(x).
\]
Because \(D(x)>0\) for \(x>1\), \(I\) decreases and then increases with the
same turning point as \(r\).  Moreover \(\lim_{x\downarrow1}I(x)=0\), so
\(I(x)<0\) immediately to the right of \(1\).

Finally,
\[
N_0(8)=(H_8-\gamma)\log5-\frac{11}{70}\log(40320)>0.
\]
For instance, the bounds
\[
\gamma<\frac{5773}{10000},\quad \log5>\frac{1609}{1000},
\quad \log(40320)<\frac{10605}{1000}
\]
give the lower bound
\[
\left(\frac{761}{280}-\frac{5773}{10000}\right)\frac{1609}{1000}
-\frac{11}{70}\frac{10605}{1000}
=\frac{124435951}{70000000}>0.
\]
Thus \(I\), and therefore \(N_0\) and \(N\), has exactly one zero on
\((1,8]\).  Denote it by \(\xi\).  The sign pattern is
\[
N(x)<0\quad(1<x<\xi),\qquad
N(\xi)=0,\qquad
N(x)>0\quad(\xi<x\le8).
\]

A high-precision numerical check places
\[
\xi=1.12620706105164346558\ldots,
\]
matching the source's reported value.  The numerical value is not used in the
proof.

The raw numerator has a removable zero at \(x=1\), since \(G(1)=D(1)=0\).  Its
quadratic leading coefficient is
\[
\lim_{x\downarrow1}\frac{N_0(x)}{(x-1)^2}
=
\frac{7(\pi^2/6-1)-11(1-\gamma)}{98}<0,
\]
again by \(\pi^2<987/100\) and \(\gamma<5773/10000\), which gives
\(-1347/980000<0\).  Therefore any normalization dividing out the removable
\((x-1)^2\) factor is negative at the left endpoint.

The single-crossing candidate is true:
\[
T\text{-BZ-gamma-quotient-N-single-crossing-one-eight}.
\]
Together with the source's theorem on \((-1,1)\) and the previous local
right-tail theorem for \(x\ge8\), this proves the critical-window reduction
\[
T\text{-BZ-gamma-quotient-critical-window-reduction}
\]
and then the source problem
\[
T\text{-Bulboaca-Zayed-gamma-quotient-full-monotonicity}.
\]

The finite-envelope candidate is not separately promoted: the proof above is a
clean analytic ratio-kernel argument rather than the finite rational cover
specified in that candidate.  The diagnostic obstruction-map candidate is also
not needed after the single-crossing proof.
\end{proof}

\begin{proposition}
\label{res:q2-analytic-denominator-reduced-assembly-terminal-certificate}
The terminal endpoint follows from a reduced analytic-denominator assembly: the true covers \((0,9/50]\) and \([9/10,1)\), the endpoint signs on \(J\), an analytic-denominator or direct-log finite-middle cover of \([9/50,1409/5000]\cup[293/1000,9/10]\), and a compact one-crossing/dominance certificate on \([1409/5000,293/1000]\setminus J\). This determines \(L_2=Q_2(\xi)\) and \(\mathcal I_2=(Q_2(\xi),3]\).
\end{proposition}
\begin{proof}
This roll proves a true auxiliary finite-middle outside-cover node:

the corresponding result.

The replayable certificate is:

\begin{verbatim}
\end{verbatim}

It proves
\[
Q_2(x)<q_J
\]
on
\[
\left[\frac9{50},\frac{1409}{5000}\right]
\cup
\left[\frac{293}{1000},\frac9{10}\right].
\]

The left bridge \([9/50,221050/1000000]\) uses the already-derived identity
\[
3xZ_4(x)-Z_3(x)
=2x^{-3}+\sum_{k=1}^{\infty}\frac{2x-k}{(x+k)^4}.
\]
For \(x<1/2\), this gives
\[
R(x)>Z_3(1/x).
\]
The certificate proves
\[
Z_3(1/x)>x^{1157/500}
\]
on \([9/50,221050/1000000]\). Since
\[
\frac{1157}{500}<\frac{115727}{50000}<q_J,
\]
this gives \(Q_2(x)<q_J\) on the bridge.

The remaining two subintervals use the direct comparison
\[
\log R(x)>\frac{115727}{50000}\log x.
\]
Since \(\log x<0\), this implies
\[
Q_2(x)<\frac{115727}{50000}<q_J.
\]
The script uses power-of-two range reduction for the atanh logarithm enclosure:
\[
\log r=\log(2^k r)-k\log 2,
\]
so every logarithm is evaluated with an argument in \([1,2)\).

The certificate output was:

\begin{verbatim}
REMAINING_FINITE_MIDDLE_COVER_OK
ZETA_LEFT_SUBCOVER
interval_count 245
domain [180000/1000000,221050/1000000]
min_width 1 max_width 1283
worst_interval [180000/1000000,181282/1000000]
worst_margin_positive True
DIRECT_LEFT_SUBCOVER
interval_count 3739
domain [221050/1000000,281800/1000000]
min_width 1 max_width 238
worst_interval [280389/1000000,280392/1000000]
worst_margin_positive True
DIRECT_RIGHT_SUBCOVER
interval_count 3801
domain [293000/1000000,900000/1000000]
min_width 1 max_width 18969
worst_interval [293124/1000000,293126/1000000]
worst_margin_positive True
\end{verbatim}

Together with the earlier true covers \((0,9/50]\) and \([9/10,1)\), this covers all outside-compact points below \(q_J\).

This roll also proves:

the corresponding result.

The replayable certificate is:

\begin{verbatim}
\end{verbatim}

Let
\[
H(x)=\Lambda(x)+x\Lambda'(x).
\]
Using \(Z_s'(x)=-sZ_{s+1}(x)\) and
\[
\frac{d}{dx}Z_s(1/x)=s x^{-2} Z_{s+1}(1/x),
\]
the script builds a rational interval enclosure for \(H\) using \(Z_3,\ldots,Z_6\) at \(x\) and \(Z_3,\ldots,Z_5\) at \(1/x\). It proves
\[
H(x)>0
\qquad
\left(\frac{1409}{5000}\le x\le\frac{293}{1000}\right).
\]

Since
\[
G'(x)=\log x\,H(x)
\]
and \(\log x<0\) on the compact bracket, this gives
\[
G'(x)<0.
\]
The existing endpoint certificate proves
\[
G(287345/1000000)>0,
\qquad
G(287346/1000000)<0.
\]
Therefore \(G\) has a unique zero
\[
\xi\in J=\left[\frac{287345}{1000000},\frac{287346}{1000000}\right],
\]
with \(G>0\) to the left of \(J\) and \(G<0\) to the right of \(J\) on the compact bracket.

Because
\[
Q_2'(x)=\frac{G(x)}{x(\log x)^2},
\]
\(Q_2\) is increasing before \(\xi\) and decreasing after \(\xi\) on the compact bracket.

The certificate output was:

\begin{verbatim}
COMPACT_MONOTONICITY_CERTIFICATE_OK
endpoint_signs G_left_positive G_right_negative
h_interval_count 128
domain [281800/1000000,293000/1000000]
min_width 87 max_width 88
worst_interval [281887/1000000,281975/1000000]
worst_H_lower_positive True
\end{verbatim}

The true cover package is now:

\[
(0,9/50]\cup
\left[9/50,\frac{1409}{5000}\right]\cup
\left[\frac{293}{1000},9/10\right]\cup
[9/10,1).
\]

All points outside the compact bracket have \(Q_2(x)<q_J\), and \(q_J<Q_2(\xi)\) by the certified inner witness. On the compact bracket, \(Q_2\) has a unique maximum at the unique zero \(\xi\in J\). Thus
\[
L_2=Q_2(\xi).
\]

At \(\beta=Q_2(\xi)\), the defining inequality has equality at \(x=\xi\), so the lower endpoint is excluded. The known upper endpoint remains included. Therefore
\[
\mathcal I_2=(Q_2(\xi),3].
\]

This proves the terminal exact endpoint node.
\end{proof}

\begin{proposition}
\label{res:q2-compact-taylor-monotonicity-j-terminal-certificate}
Using the true endpoint-sign certificate on \(J=[287345/1000000,287346/1000000]\), there is a Taylor-model or derivative-bounded rational certificate proving \(G>0\) on \([1409/5000,287345/1000000]\), \(G<0\) on \([287346/1000000,293/1000]\), and dominance below the witness \(q_J\); combined with finite-middle covers outside the compact bracket and the true endpoint outside covers, this determines \(L_2=Q_2(\xi)\) and \(\mathcal I_2=(Q_2(\xi),3]\).
\end{proposition}
\begin{proof}
This roll proves a true auxiliary finite-middle outside-cover node:

the corresponding result.

The replayable certificate is:

\begin{verbatim}
\end{verbatim}

It proves
\[
Q_2(x)<q_J
\]
on
\[
\left[\frac9{50},\frac{1409}{5000}\right]
\cup
\left[\frac{293}{1000},\frac9{10}\right].
\]

The left bridge \([9/50,221050/1000000]\) uses the already-derived identity
\[
3xZ_4(x)-Z_3(x)
=2x^{-3}+\sum_{k=1}^{\infty}\frac{2x-k}{(x+k)^4}.
\]
For \(x<1/2\), this gives
\[
R(x)>Z_3(1/x).
\]
The certificate proves
\[
Z_3(1/x)>x^{1157/500}
\]
on \([9/50,221050/1000000]\). Since
\[
\frac{1157}{500}<\frac{115727}{50000}<q_J,
\]
this gives \(Q_2(x)<q_J\) on the bridge.

The remaining two subintervals use the direct comparison
\[
\log R(x)>\frac{115727}{50000}\log x.
\]
Since \(\log x<0\), this implies
\[
Q_2(x)<\frac{115727}{50000}<q_J.
\]
The script uses power-of-two range reduction for the atanh logarithm enclosure:
\[
\log r=\log(2^k r)-k\log 2,
\]
so every logarithm is evaluated with an argument in \([1,2)\).

The certificate output was:

\begin{verbatim}
REMAINING_FINITE_MIDDLE_COVER_OK
ZETA_LEFT_SUBCOVER
interval_count 245
domain [180000/1000000,221050/1000000]
min_width 1 max_width 1283
worst_interval [180000/1000000,181282/1000000]
worst_margin_positive True
DIRECT_LEFT_SUBCOVER
interval_count 3739
domain [221050/1000000,281800/1000000]
min_width 1 max_width 238
worst_interval [280389/1000000,280392/1000000]
worst_margin_positive True
DIRECT_RIGHT_SUBCOVER
interval_count 3801
domain [293000/1000000,900000/1000000]
min_width 1 max_width 18969
worst_interval [293124/1000000,293126/1000000]
worst_margin_positive True
\end{verbatim}

Together with the earlier true covers \((0,9/50]\) and \([9/10,1)\), this covers all outside-compact points below \(q_J\).

This roll also proves:

the corresponding result.

The replayable certificate is:

\begin{verbatim}
\end{verbatim}

Let
\[
H(x)=\Lambda(x)+x\Lambda'(x).
\]
Using \(Z_s'(x)=-sZ_{s+1}(x)\) and
\[
\frac{d}{dx}Z_s(1/x)=s x^{-2} Z_{s+1}(1/x),
\]
the script builds a rational interval enclosure for \(H\) using \(Z_3,\ldots,Z_6\) at \(x\) and \(Z_3,\ldots,Z_5\) at \(1/x\). It proves
\[
H(x)>0
\qquad
\left(\frac{1409}{5000}\le x\le\frac{293}{1000}\right).
\]

Since
\[
G'(x)=\log x\,H(x)
\]
and \(\log x<0\) on the compact bracket, this gives
\[
G'(x)<0.
\]
The existing endpoint certificate proves
\[
G(287345/1000000)>0,
\qquad
G(287346/1000000)<0.
\]
Therefore \(G\) has a unique zero
\[
\xi\in J=\left[\frac{287345}{1000000},\frac{287346}{1000000}\right],
\]
with \(G>0\) to the left of \(J\) and \(G<0\) to the right of \(J\) on the compact bracket.

Because
\[
Q_2'(x)=\frac{G(x)}{x(\log x)^2},
\]
\(Q_2\) is increasing before \(\xi\) and decreasing after \(\xi\) on the compact bracket.

The certificate output was:

\begin{verbatim}
COMPACT_MONOTONICITY_CERTIFICATE_OK
endpoint_signs G_left_positive G_right_negative
h_interval_count 128
domain [281800/1000000,293000/1000000]
min_width 87 max_width 88
worst_interval [281887/1000000,281975/1000000]
worst_H_lower_positive True
\end{verbatim}

The true cover package is now:

\[
(0,9/50]\cup
\left[9/50,\frac{1409}{5000}\right]\cup
\left[\frac{293}{1000},9/10\right]\cup
[9/10,1).
\]

All points outside the compact bracket have \(Q_2(x)<q_J\), and \(q_J<Q_2(\xi)\) by the certified inner witness. On the compact bracket, \(Q_2\) has a unique maximum at the unique zero \(\xi\in J\). Thus
\[
L_2=Q_2(\xi).
\]

At \(\beta=Q_2(\xi)\), the defining inequality has equality at \(x=\xi\), so the lower endpoint is excluded. The known upper endpoint remains included. Therefore
\[
\mathcal I_2=(Q_2(\xi),3].
\]

This proves the terminal exact endpoint node.
\end{proof}

\begin{proposition}
\label{res:q2-remaining-finite-middle-log-cover-terminal-certificate}
Using the true outside covers \((0,9/50]\) and \([9/10,1)\), the true endpoint signs on \(J\), and the true enclosure lemmas, there is a finite rational direct-log cover of \([9/50,1409/5000]\cup[293/1000,9/10]\) proving \(Q_2(x)<q_J\) on every interval, together with a compact one-crossing/dominance certificate on \([1409/5000,293/1000]\setminus J\). This determines \(L_2=Q_2(\xi)\) and \(\mathcal I_2=(Q_2(\xi),3]\).
\end{proposition}
\begin{proof}
This roll proves a true auxiliary finite-middle outside-cover node:

the corresponding result.

The replayable certificate is:

\begin{verbatim}
\end{verbatim}

It proves
\[
Q_2(x)<q_J
\]
on
\[
\left[\frac9{50},\frac{1409}{5000}\right]
\cup
\left[\frac{293}{1000},\frac9{10}\right].
\]

The left bridge \([9/50,221050/1000000]\) uses the already-derived identity
\[
3xZ_4(x)-Z_3(x)
=2x^{-3}+\sum_{k=1}^{\infty}\frac{2x-k}{(x+k)^4}.
\]
For \(x<1/2\), this gives
\[
R(x)>Z_3(1/x).
\]
The certificate proves
\[
Z_3(1/x)>x^{1157/500}
\]
on \([9/50,221050/1000000]\). Since
\[
\frac{1157}{500}<\frac{115727}{50000}<q_J,
\]
this gives \(Q_2(x)<q_J\) on the bridge.

The remaining two subintervals use the direct comparison
\[
\log R(x)>\frac{115727}{50000}\log x.
\]
Since \(\log x<0\), this implies
\[
Q_2(x)<\frac{115727}{50000}<q_J.
\]
The script uses power-of-two range reduction for the atanh logarithm enclosure:
\[
\log r=\log(2^k r)-k\log 2,
\]
so every logarithm is evaluated with an argument in \([1,2)\).

The certificate output was:

\begin{verbatim}
REMAINING_FINITE_MIDDLE_COVER_OK
ZETA_LEFT_SUBCOVER
interval_count 245
domain [180000/1000000,221050/1000000]
min_width 1 max_width 1283
worst_interval [180000/1000000,181282/1000000]
worst_margin_positive True
DIRECT_LEFT_SUBCOVER
interval_count 3739
domain [221050/1000000,281800/1000000]
min_width 1 max_width 238
worst_interval [280389/1000000,280392/1000000]
worst_margin_positive True
DIRECT_RIGHT_SUBCOVER
interval_count 3801
domain [293000/1000000,900000/1000000]
min_width 1 max_width 18969
worst_interval [293124/1000000,293126/1000000]
worst_margin_positive True
\end{verbatim}

Together with the earlier true covers \((0,9/50]\) and \([9/10,1)\), this covers all outside-compact points below \(q_J\).

This roll also proves:

the corresponding result.

The replayable certificate is:

\begin{verbatim}
\end{verbatim}

Let
\[
H(x)=\Lambda(x)+x\Lambda'(x).
\]
Using \(Z_s'(x)=-sZ_{s+1}(x)\) and
\[
\frac{d}{dx}Z_s(1/x)=s x^{-2} Z_{s+1}(1/x),
\]
the script builds a rational interval enclosure for \(H\) using \(Z_3,\ldots,Z_6\) at \(x\) and \(Z_3,\ldots,Z_5\) at \(1/x\). It proves
\[
H(x)>0
\qquad
\left(\frac{1409}{5000}\le x\le\frac{293}{1000}\right).
\]

Since
\[
G'(x)=\log x\,H(x)
\]
and \(\log x<0\) on the compact bracket, this gives
\[
G'(x)<0.
\]
The existing endpoint certificate proves
\[
G(287345/1000000)>0,
\qquad
G(287346/1000000)<0.
\]
Therefore \(G\) has a unique zero
\[
\xi\in J=\left[\frac{287345}{1000000},\frac{287346}{1000000}\right],
\]
with \(G>0\) to the left of \(J\) and \(G<0\) to the right of \(J\) on the compact bracket.

Because
\[
Q_2'(x)=\frac{G(x)}{x(\log x)^2},
\]
\(Q_2\) is increasing before \(\xi\) and decreasing after \(\xi\) on the compact bracket.

The certificate output was:

\begin{verbatim}
COMPACT_MONOTONICITY_CERTIFICATE_OK
endpoint_signs G_left_positive G_right_negative
h_interval_count 128
domain [281800/1000000,293000/1000000]
min_width 87 max_width 88
worst_interval [281887/1000000,281975/1000000]
worst_H_lower_positive True
\end{verbatim}

The true cover package is now:

\[
(0,9/50]\cup
\left[9/50,\frac{1409}{5000}\right]\cup
\left[\frac{293}{1000},9/10\right]\cup
[9/10,1).
\]

All points outside the compact bracket have \(Q_2(x)<q_J\), and \(q_J<Q_2(\xi)\) by the certified inner witness. On the compact bracket, \(Q_2\) has a unique maximum at the unique zero \(\xi\in J\). Thus
\[
L_2=Q_2(\xi).
\]

At \(\beta=Q_2(\xi)\), the defining inequality has equality at \(x=\xi\), so the lower endpoint is excluded. The known upper endpoint remains included. Therefore
\[
\mathcal I_2=(Q_2(\xi),3].
\]

This proves the terminal exact endpoint node.
\end{proof}

\begin{proposition}
\label{res:zeta-tail-s9-reusable-certificate}
The next case \(s=9\) has an exact formula for \(\lfloor\zeta_n(9)^{-1}\rfloor\) with a reproducible certificate following the same reusable approximant, telescoping-sign, and gap-trap template intended for later exponents.
\end{proposition}
\begin{proof}
For integer \(s>1\), the Hurwitz-tail Euler--Maclaurin expansion has the form
\[
\zeta_n(s)
=n^{1-s}\left[
\frac1{s-1}+\frac{z}{2}
+\sum_{r\ge1}\frac{B_{2r}}{(2r)!}(s)_{2r-1}z^{2r}
\right],
\qquad z=\frac1n.
\]
Formally inverting the bracket gives
\[
\zeta_n(s)^{-1}\sim n^{s-1}G_s(z).
\]

The replay script
computes \(G_s\) symbolically and verifies two things:

1. For \(s=8\), the coefficients through \(z^{15}\) reproduce the staged \(A_8\) and \(B_8-A_8\) coefficients.
2. For \(s=9\), truncating \(n^8G_9(1/n)\) through \(z^{14}\) gives an \(A_9\), and adding the \(z^{15}\) term gives \(B_9\), with exact telescoping residuals of the expected signs.

The generated \(s=9\) lower approximant is
\[
\begin{aligned}
A_9(n)=&8n^8-32n^7+80n^6-128n^5+120n^4-64n^3
+\frac{624}{7}n^2-\frac{512}{7}n-\frac{3324}{7}\\
&+\frac{90304}{21n^2}+\frac{90304}{21n^3}
-\frac{10280944}{245n^4}-\frac{64846304}{735n^5}
+\frac{109940976}{245n^6}.
\end{aligned}
\]
The upper approximant is
\[
B_9(n)=A_9(n)+\frac{384388288}{245n^7}.
\]

Let \(R_A(n)=1/A_9(n)\) and \(R_B(n)=1/B_9(n)\). Exact denominator clearing gives:
\[
R_A(k)-R_A(k+1)-\frac1{k^9}>0
\qquad(k\ge5),
\]
because the numerator, after substituting \(k=m+5\), has all coefficients positive. Similarly,
\[
\frac1{k^9}-\left(R_B(k)-R_B(k+1)\right)>0
\qquad(k\ge1),
\]
because the denominator-cleared numerator, after substituting \(k=m+1\), has all coefficients positive.

Therefore, for \(n\ge5\),
\[
A_9(n)<\zeta_n(9)^{-1}<B_9(n).
\]

This proves the inverse-window part of the reusable template and promotes the corresponding result.

The replay script also verifies the no-integer gap.

Write \(P_9(n)\) for the polynomial part of \(A_9(n)\):
\[
P_9(n)=8n^8-32n^7+80n^6-128n^5+120n^4-64n^3
+\frac{624}{7}n^2-\frac{512}{7}n-\frac{3324}{7}.
\]
Modulo \(7\), the fractional numerator of \(P_9(n)\) is
\[
n^2-n+1\pmod 7,
\]
so the fractional part of \(P_9(n)\) is one of
\[
0,\quad \frac17,\quad \frac37,\quad \frac67.
\]

The correction \(A_9(n)-P_9(n)\) is positive for \(n\ge9\); after denominator clearing and substituting \(n=m+9\), the numerator has all coefficients positive. For \(n\ge174\), the upper correction is bounded by
\[
B_9(n)-P_9(n)
<
\frac{90304}{21n^2}
+\frac{90304}{21n^3}
+\frac{109940976}{245n^6}
+\frac{384388288}{245n^7}
<\frac17.
\]
Therefore \(A_9(n)\) and \(B_9(n)\) cannot straddle an integer for \(n\ge174\). Exact rational arithmetic verifies
\[
\lfloor A_9(n)\rfloor=\lfloor B_9(n)\rfloor
\qquad(9\le n\le173).
\]
Consequently,
\[
\left\lfloor\zeta_n(9)^{-1}\right\rfloor=\lfloor A_9(n)\rfloor
\qquad(n\ge9).
\]

For \(1\le n\le8\), the exact values are
\[
\begin{array}{c|rrrrrrrr}
n&1&2&3&4&5&6&7&8\\
\hline
\lfloor\zeta_n(9)^{-1}\rfloor&0&497&18093&224086&1543530&7360226&27295287&84349541.
\end{array}
\]
For \(n=1\), this follows from \(\zeta(9)>1\). For \(2\le n\le8\), the replay script uses exact rational partial sums through \(N=100\) and the integral tail bound
\[
\sum_{k=N+1}^{\infty}k^{-9}<\frac{1}{8N^8}
\]
to prove
\[
\frac1{M_n+1}<\zeta_n(9)<\frac1{M_n}.
\]

This proves the claim.
\end{proof}

\begin{proposition}
\label{res:yang-detemple-r-n0-cm-certificate}
The function \(F_0(x)=\left(24x^2+21/5\right)(\psi(x+1/2)-\log x)-1\) is completely monotone on \((0,\infty)\), with nonzero finite tail coefficient \(2071/33600\) at order \(x^{-4}\).
\end{proposition}
\begin{proof}
The standard asymptotic expansion of \(\psi\) gives
\[
R(x)=\frac{1}{24x^2}-\frac{7}{960x^4}+O(x^{-6}).
\]
For the normalized \(n=0\) expression
\[
F(x)=(a_0+24x^2)R(x)-b_0
\]
to have a finite nonzero \(x^{-4}\) first tail term, the constant and \(x^{-2}\) terms must vanish.  Hence
\[
b_0=1,
\]
and
\[
\frac{a_0}{24}+24\left(-\frac{7}{960}\right)=0,
\]
so
\[
a_0=\frac{21}{5}.
\]
Thus the normalized candidate is
\[
F_0(x)=\left(24x^2+\frac{21}{5}\right)R(x)-1.
\]

Using
\[
\psi(z)=\int_0^\infty\left(\frac{e^{-t}}{t}-\frac{e^{-zt}}{1-e^{-t}}\right)dt
\]
and
\[
\log x=\int_0^\infty\frac{e^{-t}-e^{-xt}}{t}\,dt,
\]
we get
\[
R(x)=\int_0^\infty e^{-xt}r(t)\,dt,
\qquad
r(t)=\frac1t-\frac{1}{2\sinh(t/2)}.
\]
The kernel has \(r(0)=0\) and \(r'(0)=1/24\).  Twice integrating by parts gives
\[
x^2R(x)-\frac1{24}=\int_0^\infty e^{-xt}r''(t)\,dt.
\]
Therefore
\[
F_0(x)=\int_0^\infty e^{-xt}K(t)\,dt,
\qquad
K(t)=24r''(t)+\frac{21}{5}r(t).
\]

It remains to prove \(K(t)\ge0\).

Put \(y=t/2\) and \(h(y)=1/y-\operatorname{csch}y\).  Since \(r(t)=h(y)/2\) and \(d/dt=(1/2)d/dy\),
\[
K(t)=3h''(y)+\frac{21}{10}h(y).
\]
With
\[
h''(y)=\frac{2}{y^3}-\operatorname{csch}y\,\coth^2y-\operatorname{csch}^3y
\]
and \(\coth^2y=1+\operatorname{csch}^2y\), this is equivalent to
\[
K(t)=
\frac{
S(y)
}{
10y^3\sinh^3y
},
\]
where
\[
S(y)=(60+21y^2)\sinh^3y-51y^3\sinh^2y-60y^3.
\]

Using
\[
\sinh^3y=\frac{\sinh3y-3\sinh y}{4},
\qquad
\sinh^2y=\frac{\cosh2y-1}{2},
\]
one obtains the power series
\[
S(y)=\sum_{m=4}^{\infty} c_my^{2m+1},
\]
where
\[
c_m=
\frac{
15(3^{2m+1}-3)
+\frac{21}{4}(3^{2m-1}-3)(2m)(2m+1)
-\frac{51}{2}2^{2m-2}(2m-1)(2m)(2m+1)
}{
(2m+1)!}.
\]
The first coefficients are already positive:
\[
c_4=\frac{2071}{2520},\qquad
c_5=\frac{1909}{10080}.
\]
For \(m\ge6\), the positive \(3^{2m-1}\)-terms dominate the \(2^{2m-2}\) term; explicitly, after multiplying by \((2m+1)!\), the ratio of the positive part to the negative part is increasing from \(m=6\) onward and is already \(>1\) at \(m=6\).  Thus \(c_m>0\) for every \(m\ge4\).

Hence \(S(y)>0\) for \(y>0\), and therefore \(K(t)>0\) for \(t>0\).

By Bernstein's theorem,
\[
F_0(x)=\left(24x^2+\frac{21}{5}\right)(\psi(x+1/2)-\log x)-1
\]
is completely monotone on \((0,\infty)\).

Since \(K(t)=\frac{2071}{201600}t^3+O(t^5)\), Watson's lemma gives
\[
F_0(x)=\frac{2071}{33600}x^{-4}+O(x^{-6}).
\]
This verifies the nonzero finite tail condition for \(n=0\).
\end{proof}

\begin{proposition}
\label{res:cbf-difference-quotient-stieltjes}
If \(G\) is a complete Bernstein function with finite \(G(0+)\), then \((G(x)-G(0+))/x\) is a Stieltjes function.
\end{proposition}
\begin{proof}
Let \(q\ge1\) and \(b_i>a_i>0\). Define
\[
F(x)={}_{q+1}F_q(1,a_1,\ldots,a_q;b_1,\ldots,b_q;-x)
\]
and
\[
R(x)=
\frac{{}_{q+1}F_q(1,a_1+1,\ldots,a_q+1;b_1+1,\ldots,b_q+1;-x)}
{{}_{q+1}F_q(1,a_1,\ldots,a_q;b_1,\ldots,b_q;-x)}.
\]
Then \(R\) is a Stieltjes function on \((0,\infty)\), hence completely monotone. Moreover
\[
xR(x)
\]
is a complete Bernstein function up to the positive scalar \(\prod_i b_i/a_i\).

For \(q=1\), this gives the Gauss/Beta slice:
\[
\frac{{}_2F_1(1,a+1;b+1;-x)}{{}_2F_1(1,a;b;-x)}
\]
is Stieltjes and completely monotone for \(b>a>0\).

Karp--Sitnik's positive-density representation gives a compact-support Stieltjes form for \(F\):
\[
F(x)=\int_0^1 \frac{d\mu(s)}{1+sx},
\qquad \mu\ge0,
\qquad \mu([0,1])=F(0)=1.
\]
Thus \(F\) is a nonzero Stieltjes function.

By the standard Stieltjes/complete-Bernstein duality, if \(F\) is nonzero Stieltjes, then
\[
G(x)=\frac1{F(x)}
\]
is a complete Bernstein function. Since \(G(0)=1\), the complete Bernstein representation gives
\[
G(x)=1+\alpha x+\int_0^\infty \frac{x}{x+t}\,d\nu(t),
\qquad \alpha\ge0,\quad \nu\ge0.
\]
Therefore
\[
\frac{G(x)-1}{x}
=\alpha+\int_0^\infty \frac{1}{x+t}\,d\nu(t)
\]
is Stieltjes.

Karp--Sitnik's contiguous defect identity for \(\sigma=1,\delta=1\) is
\[
1-F(x)
=x\,{}_{q+1}F_q(1,a_1+1,\ldots,a_q+1;b_1+1,\ldots,b_q+1;-x)
\prod_{i=1}^q\frac{a_i}{b_i}.
\]
Hence
\[
R(x)=
\prod_{i=1}^q\frac{b_i}{a_i}\,
\frac{1-F(x)}{xF(x)}
=
\prod_{i=1}^q\frac{b_i}{a_i}\,
\frac{G(x)-1}{x}.
\]
A positive scalar multiple of a Stieltjes function is Stieltjes. Thus \(R\) is Stieltjes and completely monotone.

Finally,
\[
xR(x)=\prod_{i=1}^q\frac{b_i}{a_i}\,(G(x)-1).
\]
Since subtracting the finite value \(G(0)=1\) from a complete Bernstein function leaves the same nonnegative drift and Levy measure part, \(G(x)-1\) is complete Bernstein. Hence \(xR(x)\) is complete Bernstein up to the stated positive scalar.

This is a strong local theorem and a reusable Stieltjes/CBF bridge. The subsequent primary-source wording audit did not find an explicit Karp--Sitnik open problem asking for the CM/Stieltjes/CBF upgrade, so this result is bridge-only unless a separate source asks for exactly this strengthening.
\end{proof}

\begin{proposition}
\label{res:baskakov-r4-alpha1-laplace-density-seed}
For Abel--Gawronski--Neuschel's higher-power Baskakov function, \(f^{[4]}_1(x)=1/((1+x)^4-x^4)=\int_0^\infty e^{-xt}e^{-t/2}(1-\cos(t/2))\,dt\); hence \(f^{[4]}_1\) is completely monotone on \((0,\infty)\).
\end{proposition}
\begin{proof}
For \(\alpha=1\),
\[
\binom{-1}{k}=(-1)^k,
\qquad
\binom{-1}{k}^{4}=1.
\]
Therefore, for \(x>0\),
\[
f^{[4]}_1(x)
=(1+x)^{-4}\sum_{k=0}^{\infty}
\left(\frac{x}{1+x}\right)^{4k}
=\frac{1}{(1+x)^4-x^4}.
\]
Since
\[
(1+x)^4-x^4=(2x+1)(2x^2+2x+1),
\]
partial fractions give
\[
\frac{1}{(1+x)^4-x^4}
=\frac{2}{2x+1}-\frac{2x+1}{2x^2+2x+1}.
\]
Writing \(a=x+\frac12\), this is
\[
\frac{1}{a}-\frac{a}{a^2+(1/2)^2}.
\]
Using
\[
\int_0^\infty e^{-xt}e^{-t/2}\,dt=\frac{1}{x+1/2}
\]
and
\[
\int_0^\infty e^{-xt}e^{-t/2}\cos(t/2)\,dt
=\frac{x+1/2}{(x+1/2)^2+(1/2)^2},
\]
we obtain
\[
f^{[4]}_1(x)
=\int_0^\infty e^{-xt}e^{-t/2}\left(1-\cos\frac{t}{2}\right)\,dt.
\]
The density \(e^{-t/2}(1-\cos(t/2))\) is nonnegative on \([0,\infty)\). By Bernstein's theorem, \(f^{[4]}_1\) is completely monotone on \((0,\infty)\).

From the source definition
\[
p^{[c]}_{n,k}(x)=\binom{-n/c}{k}(-cx)^k(1+cx)^{-n/c-k},
\]
one has
\[
\psi^{[r]}_{n,c}(x)=\sum_{k=0}^{\infty}\left(p^{[c]}_{n,k}(x)\right)^r
=f^{[r]}_{n/c}(cx)
\]
for even \(r\). Hence, for \(c=n\),
\[
\psi^{[4]}_{n,n}(x)=f^{[4]}_1(nx).
\]
Positive scaling preserves complete monotonicity, so \(\psi^{[4]}_{n,n}\) is completely monotone for every positive integer \(n\).
\end{proof}

\begin{proposition}
\label{res:bouali-falpha-phi-asymptotic}
For \(f_\alpha(x)=x^{x(\psi(x)-\log x)-\alpha}\), \(\theta(x)=x(\log x-\psi(x))\), and \(\phi_\alpha(x)=-(\log f_\alpha)'(x)\), one has \(\phi_\alpha(x)=(\alpha+1/2)/x+(1-\log x)/(12x^2)+O(\log x/x^4)\) as \(x\to\infty\).
\end{proposition}
\begin{proof}
Status: endpoint/lower-region obstruction solved; interior gap remains open.

Bouali studies
\[
f_\alpha(x)=x^{x(\psi(x)-\log x)-\alpha}.
\]
The source proves logarithmic complete monotonicity, hence complete monotonicity, for \(\alpha\ge -1/4\). Remark 1.12 prints the open range as \((-1/4,-1/2]\), which is reversed as a literal interval. The local correction treats the intended unresolved lower side as a frontier and proves that the endpoint \(\alpha=-1/2\) cannot be included.

Let \(\theta(x)=x(\log x-\psi(x))\) and
\[
\phi_\alpha(x)=-(\log f_\alpha)'(x)=\frac{\theta(x)+\alpha}{x}+\theta'(x)\log x.
\]
The digamma expansion gives
\[
\phi_\alpha(x)=\frac{\alpha+1/2}{x}+\frac{1-\log x}{12x^2}+O\!\left(\frac{\log x}{x^4}\right).
\]
Thus \(\phi_\alpha(x)<0\) eventually whenever \(\alpha\le -1/2\). Since positive complete monotone functions must be nonincreasing, this proves
\[
\alpha\le -\frac12 \quad\Longrightarrow\quad f_\alpha\notin CM(0,\infty).
\]

This is useful theory growth because it packages a reusable Gamma/psi endpoint-obstruction template, but it is not application-eligible as a full solution. The interior interval
\[
-\frac12<\alpha<-\frac14
\]
remains open.

confirmed the corrected interior interval as source-open and found no later
gap. It isolated the reduction
\[
f_\alpha(x)=x^\beta f_{-1/4}(x),
\qquad
\beta=-(\alpha+1/4)\in(0,1/4),
\]
so the remaining problem is whether the non-CM factor \(x^\beta\) preserves
complete monotonicity for this specific endpoint function. This is useful
theory growth, but not a solved source open problem and not APP-countable.
\end{proof}

\begin{proposition}
\label{res:gt-bf-factorized-denominator-power-closure}
Let \(0<\alpha<1\), let \(\psi\in BF\), and suppose \(x/\psi(x)=\prod_{j=1}^m f_j(x)\), where every \(f_j\) is a positive Bernstein function on \((0,\infty)\). Then \(\psi_\alpha(x)=[\psi(x^\alpha)]^{1/\alpha}\) is a Bernstein function.
\end{proposition}
\begin{proof}
Assume \(0<\alpha<1\), \(\psi\in BF\), and
\[
\frac{x}{\psi(x)}=\prod_{j=1}^m f_j(x),
\]
where every \(f_j\) is a positive Bernstein function on \((0,\infty)\). Then
\[
\psi_\alpha(x)=[\psi(x^\alpha)]^{1/\alpha}
\]
is a Bernstein function.

Put
\[
\beta=\frac1\alpha-1>0.
\]
For \(x>0\), direct differentiation gives
\[
\psi_\alpha'(x)
=x^{\alpha-1}\psi'(x^\alpha)\psi(x^\alpha)^{1/\alpha-1}
=\psi'(x^\alpha)\left(\frac{\psi(x^\alpha)}{x^\alpha}\right)^\beta.
\]
Using the factorization of \(x/\psi(x)\), this becomes
\[
\psi_\alpha'(x)
=\psi'(x^\alpha)\prod_{j=1}^m f_j(x^\alpha)^{-\beta}.
\]

Since \(\psi\in BF\), \(\psi'\) is completely monotone. Since \(x^\alpha\in BF\), composition closure gives \(\psi'(x^\alpha)\in CM\).

For each \(j\), \(g_j(x)=f_j(x^\alpha)\) is Bernstein by BF composition closure. If \(g\) is a positive Bernstein function and \(\beta>0\), then
\[
g(x)^{-\beta}
=\frac1{\Gamma(\beta)}
\int_0^\infty t^{\beta-1}e^{-t g(x)}\,dt.
\]
For each \(t>0\), \(e^{-t g(x)}\) is completely monotone because \(y\mapsto e^{-ty}\) is completely monotone and \(g\in BF\). Hence \(g^{-\beta}\in CM\). Therefore every \(f_j(x^\alpha)^{-\beta}\) is completely monotone.

Finite products of completely monotone functions are completely monotone, so \(\psi_\alpha'\in CM\). Also \(\psi_\alpha\ge0\). Thus \(\psi_\alpha\) is a Bernstein function.

The full source problem remains open locally in the arbitrary non-special case, especially \(1/2<\alpha<1\), when no finite Bernstein-factor representation of \(x/\psi(x)\) is known. The next loop should rotate or look for a genuinely broader structural theorem/counterexample, not repeat this finite-factor subclass.
\end{proof}

\begin{proposition}
\label{res:kms-arrowhead-riesz-kernel-pminus2}
For \(p(x,y,z)=x(xy-z^2)\) on \(C=\{x>0,xy>z^2\}\), under pairing \(xu+yv+zw\), \(p^{-2}\) has nonnegative Riesz kernel \(K(u,v,w)=\frac{4\sqrt v}{15\pi}(u-w^2/(4v))^{5/2}\mathbf1_{\{v>0,u>w^2/(4v)\}}\).
\end{proposition}
\begin{proof}
Let
\[
p(x,y,z)=x(xy-z^2),
\qquad
C=\{(x,y,z):x>0,\ xy>z^2\}.
\]
Use the pairing \(xu+yv+zw\). Define
\[
K(u,v,w)=
\frac{4\sqrt v}{15\pi}
\left(u-\frac{w^2}{4v}\right)^{5/2}
\mathbf 1_{\{v>0,\ u>w^2/(4v)\}}.
\]
Then, for \((x,y,z)\in C\),
\[
\int_{\mathbb R^3}e^{-xu-yv-zw}K(u,v,w)\,du\,dv\,dw
=\frac{1}{x^2(xy-z^2)^2}
=p(x,y,z)^{-2}.
\]

Set
\[
s=u-\frac{w^2}{4v}.
\]
Then \(u=s+w^2/(4v)\), \(s>0\), \(v>0\), and the transform equals
\[
\frac{4}{15\pi}\int_0^\infty\int_{-\infty}^{\infty}\int_0^\infty
e^{-xs-xw^2/(4v)-yv-zw}\sqrt v\,s^{5/2}\,ds\,dw\,dv.
\]
The \(s\)-integral is
\[
\int_0^\infty e^{-xs}s^{5/2}\,ds
=\Gamma(7/2)x^{-7/2}
=\frac{15\sqrt\pi}{8}x^{-7/2}.
\]
The \(w\)-integral is
\[
\int_{-\infty}^{\infty}\exp\left(-\frac{xw^2}{4v}-zw\right)\,dw
=2\sqrt{\frac{\pi v}{x}}\exp\left(\frac{vz^2}{x}\right).
\]
Therefore the transform becomes
\[
x^{-4}\int_0^\infty v\exp\left[-v\left(y-\frac{z^2}{x}\right)\right]\,dv.
\]
Since \((x,y,z)\in C\), \(y-z^2/x>0\), and
\[
\int_0^\infty ve^{-av}\,dv=a^{-2}.
\]
Thus
\[
x^{-4}\left(y-\frac{z^2}{x}\right)^{-2}
=\frac{1}{x^2(xy-z^2)^2}
=p(x,y,z)^{-2}.
\]

Identify
\[
(x,y,z)\leftrightarrow
\begin{pmatrix}x&z\\ z&y\end{pmatrix},
\qquad
(u,v,w)\leftrightarrow
\begin{pmatrix}u&w/2\\ w/2&v\end{pmatrix}.
\]
Then \(\operatorname{tr}(AB)=xu+yv+zw\). Hence the dual cone is
\[
C^*=\{u\ge0,\ v\ge0,\ 4uv\ge w^2\}.
\]
The formula uses the coordinate \(w=2B_{12}\). In the raw symmetric-matrix coordinate \(s=B_{12}\), the density becomes
\[
\frac{8\sqrt v}{15\pi}
\left(u-\frac{s^2}{v}\right)^{5/2}
\mathbf 1_{\{v>0,\ u>s^2/v\}}.
\]
\end{proof}

\begin{counterexample}
\label{res:keady-cm-inverse-not-cm-example}
The strictly completely monotone function \(f(x)=99e^{-x}+e^{-10x}\) maps \((0,\infty)\) decreasingly onto \((0,100)\), but its inverse branch is not completely monotone on \((0,100)\).
\end{counterexample}
\begin{proof}
There exists a strictly completely monotone decreasing function whose inverse branch is not completely monotone on its natural interval. An explicit example is
\[
f(x)=99e^{-x}+e^{-10x},\qquad x>0.
\]
It maps \((0,\infty)\) decreasingly onto \((0,100)\). Its inverse branch \(g=f^{-1}:(0,100)\to(0,\infty)\) is not completely monotone.

For every \(n\ge0\),
\[
(-1)^n f^{(n)}(x)=99e^{-x}+10^n e^{-10x}>0.
\]
Hence \(f\) is strictly completely monotone, and \(f'<0\), so the inverse branch \(g=f^{-1}\) exists.

For \(y=f(x)\), inverse differentiation gives
\[
g'(y)=\frac1{f'(x)},\qquad
g''(y)=-\frac{f''(x)}{(f'(x))^3},
\]
and
\[
g'''(y)=\frac{3(f''(x))^2-f'''(x)f'(x)}{(f'(x))^5}.
\]
At \(x=0+\),
\[
f'(0+)=-109,\qquad f''(0+)=199,\qquad f'''(0+)=-1099.
\]
Therefore
\[
3(f''(0+))^2-f'''(0+)f'(0+)
=3\cdot199^2-(-1099)(-109)
=-988<0.
\]
Since \((f'(0+))^5<0\), this gives
\[
g'''(100-)>0.
\]
By continuity, \(g'''>0\) on a left-neighborhood of \(100\). Complete monotonicity of \(g\) would require \((-1)^3g'''\ge0\), i.e. \(g'''\le0\). This contradiction proves that \(g\) is not completely monotone.
\end{proof}

\begin{counterexample}
\label{res:szabo-psi-difference-alpha2-not-cm}
For every \(0<d<1\), the function \(y\mapsto y^2\left[\psi(y+d)-\psi(y)-d/y\right]\) is not completely monotone on \((0,\infty)\).
\end{counterexample}
\begin{proof}
Let \(0<d<1\) and
\[
\phi_d(y)=\psi(y+d)-\psi(y)-\frac{d}{y},\qquad y>0.
\]
Then \(y^2\phi_d(y)\) is not completely monotone on \((0,\infty)\). Equivalently, for Szabo's original variables \(d=b-a\) and \(y=x+a\), the endpoint
\[
\alpha=2
\]
is not admissible in Open Problem 1.5 for any \(0<b-a<1\).

The standard digamma difference formula gives
\[
\psi(y+d)-\psi(y)=\int_0^\infty e^{-yt}\frac{1-e^{-dt}}{1-e^{-t}}\,dt
\]
and
\[
\frac{d}{y}=\int_0^\infty e^{-yt}d\,dt.
\]
Thus
\[
\phi_d(y)=\int_0^\infty e^{-yt}k_d(t)\,dt,\qquad
k_d(t)=\frac{1-e^{-dt}}{1-e^{-t}}-d.
\]
At the origin,
\[
k_d(t)=\frac{d(1-d)}2t+O(t^2),
\]
so \(k_d(0+)=0\) and \(k_d'(0+)=d(1-d)/2>0\).

In the sense of Laplace transforms of distributions,
\[
y^2\phi_d(y)
=\frac{d(1-d)}2+\int_0^\infty e^{-yt}k_d''(t)\,dt.
\]
Indeed, integrating by parts twice gives
\[
\mathcal L(k_d'')(y)=y^2\mathcal L(k_d)(y)-k_d'(0+),
\]
because \(k_d(0+)=0\).

As \(t\to\infty\),
\[
k_d(t)=1-d-e^{-dt}+O(e^{-t}),
\]
and hence
\[
k_d''(t)=-d^2e^{-dt}+O(e^{-t}).
\]
Since \(0<d<1\), the term \(e^{-dt}\) dominates \(e^{-t}\), so \(k_d''(t)<0\) for all sufficiently large \(t\).

If \(y^2\phi_d(y)\) were completely monotone on \((0,\infty)\), the Bernstein-Widder theorem would give a nonnegative representing measure. The inverse Laplace distribution computed above has an absolutely continuous tail with negative density, contradicting uniqueness of the Laplace transform. Therefore \(y^2\phi_d(y)\) is not completely monotone.

The shift \(y=x+a\) preserves the obstruction for the original function on \((-a,\infty)\): the representing distribution is multiplied by \(e^{-at}\), which does not change the eventual negative sign.
\end{proof}

\begin{proposition}
\label{res:sqrt-two-term-symbol-concavity-loss-criterion}
For \(g(s)=c s^a+d s^b\), \(c,d>0\), \(1<a<2\), and \(0<b<a-1\), the function \(h(s)=\sqrt{g(s)}\) has \(h''(s)>0\) somewhere if and only if \((a-1)^2+(b-1)^2>1\).
\end{proposition}
\begin{proof}
Primary source: Emilia Bazhlekova and Ivan Bazhlekov, "Subordination approach to multi-term time-fractional diffusion-wave equations", arXiv:1707.09828.

The source proves the propagation-function positivity package under
\[
1<\alpha\le2,\qquad \alpha-\alpha_m\le1,
\]
by showing that
\[
\sqrt{g(s)}\in CBF,\qquad g(s)=c s^\alpha+\sum_j c_j s^{\alpha_j}.
\]
It leaves open whether, and to what extent, this gap condition can be relaxed for the actual properties
\[
w\ge0,\qquad w_t\ge0,\qquad -w_x\ge0.
\]

For the two-term wave endpoint
\[
g(s)=c s^2+d s^b,\qquad c,d>0,\quad 0<b<1,
\]
the condition cannot be relaxed. Put \(h(s)=\sqrt{g(s)}\). Then
\[
h''(s)=\frac{d}{2\sqrt c}(b-1)(b-2)s^{b-3}+O(s^{2b-5})>0
\]
for all sufficiently large \(s\).

The source gives
\[
\mathcal L\{w_t(x,\cdot)\}(s)=e^{-xh(s)}.
\]
If \(w_t(x,\cdot)\ge0\), this Laplace transform must be completely monotone. But
\[
\frac{d^2}{ds^2}e^{-xh(s)}
=e^{-xh(s)}\left(x^2h'(s)^2-xh''(s)\right),
\]
which is negative at a large \(s_0\) after choosing \(x>0\) sufficiently small. Hence \(w_t\) fails positivity for some \(x\).

\[
g(s)=c s^a+d s^b,\qquad c,d>0,\quad 1<a<2,\quad 0<b<a-1,
\]
and set \(h(s)=\sqrt{g(s)}\), \(y=(c/d)s^{a-b}\). Then the sign of \(h''\) is the sign of
\[
N_{a,b}(y)
=a(a-2)y^2
+2(a^2-ab-a+b^2-b)y
+b(b-2).
\]
The discriminant is
\[
4(a-b)^2\left((a-1)^2+(b-1)^2-1\right).
\]
Since \(a>1\) and \(b<1\), \(h''\) is positive somewhere exactly when
\[
(a-1)^2+(b-1)^2>1.
\]
In that region, \(e^{-x\sqrt g}\) fails complete monotonicity for a suitable small \(x>0\), so \(w_t\) cannot be nonnegative for all \(x,t>0\). The source examples \((a,b)=(1.9,0.5)\) and \((1.8,0.3)\) lie in this outside-disk region.

The remaining two-term inner gap
\[
1<a<2,\qquad 0<b<a-1,\qquad (a-1)^2+(b-1)^2\le1
\]
is still open locally. There the square-root symbol is concave by this test, but no Bernstein-function representation or positivity theorem has been admitted.

\[
g(s)=s^{28/25}+s^{1/50},\qquad h(s)=\sqrt{g(s)}.
\]
Then
\[
1<a<2,\qquad 0<b<a-1,\qquad a-b=\frac{11}{10}>1,
\]
and
\[
(a-1)^2+(b-1)^2
=\left(\frac3{25}\right)^2+\left(-\frac{49}{50}\right)^2
=\frac{2437}{2500}<1.
\]
Thus this is not covered by the positive-\(h''\) outside-disk obstruction.

Direct differentiation gives
\[
h^{(5)}(1)
=-\frac{5570045943\sqrt2}{320000000000}<0.
\]
Since a Bernstein function has completely monotone derivative, it must satisfy \(h^{(5)}\ge0\). Therefore \(\sqrt{s^{28/25}+s^{1/50}}\) is not a Bernstein function.

The same example gives propagation failure, not only failure of the sufficient CBF route. With
\[
F_x(s)=e^{-xh(s)}=\mathcal L\{w_t(x,\cdot)\}(s),
\]
we have
\[
F_x^{(5)}(1)=-x h^{(5)}(1)+O(x^2)>0
\]
for all sufficiently small \(x>0\). This contradicts complete monotonicity of \(F_x\), which would be necessary if \(w_t(x,\cdot)\ge0\). Hence \(w_t\)-positivity fails for this inner-gap two-term symbol.

This is still a partial source answer, not a complete classification of the inner disk. The remaining open problem is now a residual parameter-region classification for Bernstein status and propagation positivity.
\end{proof}

\begin{proposition}
\label{res:baskakov-alpha1-even-r-frontier-open}
Determine whether \(f^{[2m]}_1(x)=1/((1+x)^{2m}-x^{2m})\) is completely monotone on \((0,\infty)\) for every integer \(m\ge2\).
\end{proposition}
\begin{proof}
For \(\alpha=1\),
\[
\binom{-1}{k}=(-1)^k,
\qquad
\binom{-1}{k}^{4}=1.
\]
Therefore, for \(x>0\),
\[
f^{[4]}_1(x)
=(1+x)^{-4}\sum_{k=0}^{\infty}
\left(\frac{x}{1+x}\right)^{4k}
=\frac{1}{(1+x)^4-x^4}.
\]
Since
\[
(1+x)^4-x^4=(2x+1)(2x^2+2x+1),
\]
partial fractions give
\[
\frac{1}{(1+x)^4-x^4}
=\frac{2}{2x+1}-\frac{2x+1}{2x^2+2x+1}.
\]
Writing \(a=x+\frac12\), this is
\[
\frac{1}{a}-\frac{a}{a^2+(1/2)^2}.
\]
Using
\[
\int_0^\infty e^{-xt}e^{-t/2}\,dt=\frac{1}{x+1/2}
\]
and
\[
\int_0^\infty e^{-xt}e^{-t/2}\cos(t/2)\,dt
=\frac{x+1/2}{(x+1/2)^2+(1/2)^2},
\]
we obtain
\[
f^{[4]}_1(x)
=\int_0^\infty e^{-xt}e^{-t/2}\left(1-\cos\frac{t}{2}\right)\,dt.
\]
The density \(e^{-t/2}(1-\cos(t/2))\) is nonnegative on \([0,\infty)\). By Bernstein's theorem, \(f^{[4]}_1\) is completely monotone on \((0,\infty)\).

From the source definition
\[
p^{[c]}_{n,k}(x)=\binom{-n/c}{k}(-cx)^k(1+cx)^{-n/c-k},
\]
one has
\[
\psi^{[r]}_{n,c}(x)=\sum_{k=0}^{\infty}\left(p^{[c]}_{n,k}(x)\right)^r
=f^{[r]}_{n/c}(cx)
\]
for even \(r\). Hence, for \(c=n\),
\[
\psi^{[4]}_{n,n}(x)=f^{[4]}_1(nx).
\]
Positive scaling preserves complete monotonicity, so \(\psi^{[4]}_{n,n}\) is completely monotone for every positive integer \(n\).
\end{proof}

\begin{proposition}
\label{res:bazhlekova-wright-density-topcap-bridge-normal-form}
For the normalized Bazhlekova symbol \(h(s)=s^\alpha(1+s^{-p})^{1/2}\), with \(0<\beta=\alpha-p/2<\alpha<1\) and the principal branch in the no-cover seed range, define \(\mathcal W_{\alpha,p}(x)=-\sum_{m\ge0}\binom{1/2}{m}x^m/\Gamma(pm-\alpha)\). Then \(\mathcal L^{-1}\{h'\}(t)=t^{-\alpha}\mathcal W_{\alpha,p}(t^p)\), and the Wright top-cap limit has the same sign as \(\mathcal W_{\alpha,p}(\lambda^{-1})\). Thus normalized derivative density positivity and top-cap positivity are the same Wright-function sign problem.
\end{proposition}
\begin{proof}
Put
\[
\alpha=\frac a2,\qquad p=a-b,\qquad \beta=\frac b2=\alpha-\frac p2,
\]
so that
\[
h(s)=s^{b/2}(1+s^{a-b})^{1/2}
=s^\alpha(1+s^{-p})^{1/2}.
\]
At the two no-cover seeds,
\[
(\alpha,p)=\left(\frac34,\frac{11}{10}\right),
\qquad
(\alpha,p)=\left(\frac{11}{20},\frac{21}{20}\right),
\]
and \(0<\beta<\alpha<1\), \(1<p<2\).

Define the Wright-normalized series
\[
\mathcal W_{\alpha,p}(x)
=
-\sum_{m=0}^{\infty}
\binom{1/2}{m}\frac{x^m}{\Gamma(pm-\alpha)}.
\]
The coefficient ratio and the gamma denominator show that this series is entire in \(x\).

For \(s>1\), the binomial expansion gives
\[
h'(s)
=
\sum_{m=0}^{\infty}
\binom{1/2}{m}(\alpha-pm)s^{\alpha-pm-1}.
\]
For each \(m\), analytic continuation of the elementary Laplace formula gives
\[
\mathcal L\left\{
\frac{t^{pm-\alpha}}{\Gamma(pm-\alpha)}
\right\}(s)
=(pm-\alpha)s^{\alpha-pm-1}.
\]
Therefore
\[
\mathcal L^{-1}\{h'\}(t)
=
-\sum_{m=0}^{\infty}
\binom{1/2}{m}
\frac{t^{pm-\alpha}}{\Gamma(pm-\alpha)}
=
t^{-\alpha}\mathcal W_{\alpha,p}(t^p).
\]
For \(1<p<2\) the Wright series has subexponential growth, so the Laplace transform is defined on \(s>0\). Equality for \(s>1\) extends to \(s>0\) by analytic continuation on the principal branch.

Consequently, at the no-cover seeds, complete monotonicity of \(h'\) is equivalent to
\[
\mathcal W_{\alpha,p}(x)\ge0
\qquad(x>0),
\]
because the inverse Laplace transform is the locally integrable density above.

The previous top-cap audit used the signed polynomials \(R_n(y)=(-1)^{n-1}Q_n(y)\). Its formal top-cap function was
\[
W^{\mathrm{top}}_{\alpha,p}(\lambda)
=
1-\frac{\Gamma(1-\alpha)}{\alpha}
\sum_{m\ge1}\binom{1/2}{m}
\frac{\lambda^{-m}}{\Gamma(pm-\alpha)}.
\]
Since
\[
-\frac1{\Gamma(-\alpha)}
=\frac{\alpha}{\Gamma(1-\alpha)},
\]
the two normalizations satisfy
\[
W^{\mathrm{top}}_{\alpha,p}(\lambda)
=
\frac{\Gamma(1-\alpha)}{\alpha}
\mathcal W_{\alpha,p}(\lambda^{-1}).
\]
The factor \(\Gamma(1-\alpha)/\alpha\) is positive. Thus top-cap positivity and normalized derivative density positivity are the same sign problem after \(x=\lambda^{-1}\).
\end{proof}

The following appendix result records the fixed source problem \textbf{APP-0037} without interrupting the main exposition.
\begin{corollary}[APP-0037: Du-Wang \(h\_3\) is increasing exactly on \(\frac12\le a\le1\)]
\label{res:du-wang-h3-open-problem-2-classification}
For \(0<a<2\), Du--Wang's function \(h_3\) is increasing on \((0,\infty)\) exactly for \(1/2\le a\le1\), and is not monotone for \(0<a<1/2\) or \(1<a<2\).

This resolves the fixed application label \textbf{APP-0037}: Classify Du-Wang \(h\_3\) monotonicity on \((0,\infty)\).
\emph{Source reference.} \cite{app-app-0037-source}.
\end{corollary}
\begin{proof}
For \(t>1/2\), define
\[
S_m(t)=\sum_{n=0}^{\infty}\frac1{(n+t)^m}.
\]
The polygamma identities give
\[
\psi''(t)=-2S_3(t),\qquad \psi'''(t)=6S_4(t).
\]
Thus the desired bound
\[
-\frac{2\psi''(t)}{\psi'''(t)}\ge t-\frac12
\]
is equivalent to
\[
2S_3(t)-3\left(t-\frac12\right)S_4(t)\ge0.
\]
Write \(y=t-\frac12>0\).  Using
\[
S_m(t)=\frac1{\Gamma(m)}
\int_0^\infty \frac{x^{m-1}e^{-tx}}{1-e^{-x}}\,dx
\]
and
\[
\frac{e^{-tx}}{1-e^{-x}}
=e^{-yx}\frac1{2\sinh(x/2)},
\]
set
\[
K(x)=\frac1{2\sinh(x/2)}.
\]
Then
\[
2S_3(t)-3yS_4(t)
=
\int_0^\infty e^{-yx}x^2K(x)\,dx
-\frac y2\int_0^\infty e^{-yx}x^3K(x)\,dx.
\]
Since \(x^3K(x)\sim x^2\) as \(x\downarrow0\) and decays exponentially at
infinity, integration by parts gives
\[
y\int_0^\infty e^{-yx}x^3K(x)\,dx
=
\int_0^\infty e^{-yx}(x^3K(x))'\,dx.
\]
Therefore
\[
2S_3(t)-3yS_4(t)
=
\int_0^\infty e^{-yx}
\left(x^2K(x)-\frac12(x^3K(x))'\right)dx.
\]
Now
\[
K'(x)=-\frac12K(x)\coth(x/2),
\]
so
\[
x^2K(x)-\frac12(x^3K(x))'
=
\frac12x^2K(x)\left(\frac x2\coth(x/2)-1\right).
\]
For \(z>0\), \(z\coth z>1\).  The integrand is therefore strictly positive,
and
\[
2S_3(t)-3\left(t-\frac12\right)S_4(t)>0.
\]
This proves
\[
-\frac{2\psi''(t)}{\psi'''(t)}>t-\frac12,
\qquad t>\frac12.
\]

This proves the claim.

Let \(1/2\le a\le1\), \(u>0\), and \(t=a+u\).  Since \(a\ge1/2\),
\[
u=t-a\le t-\frac12.
\]
The ratio lemma gives
\[
t-\frac12<-\frac{2\psi''(t)}{\psi'''(t)}.
\]
Because \(\psi'''(t)>0\),
\[
2\psi''(t)+u\psi'''(t)<0.
\]
Consequently
\[
H_a'(u)=-u^2\{2\psi''(t)+u\psi'''(t)\}>0.
\]
Together with \(H_a(0+)=2\log\Gamma(a)\ge0\), this gives
\[
H_a(u)\ge0
\qquad (u>0,\ 1/2\le a\le1).
\]

This proves the claim.

For \(x>0\), \(u=x\), and \(t=x+a\),
\[
h_3'(x)=\widetilde h_3'(x+a)=\frac{h_{31}(x+a)}{x^2}
=\frac{H_a(x)}{x^2}\ge0.
\]
The inequality is strict for \(x>0\) except possibly at the endpoint
\(a=1,u=0\), which is outside the open \(x\)-interval.  Thus \(h_3\) is
increasing on \((0,\infty)\) for every
\[
\frac12\le a\le1.
\]

This proves the claim.
the corresponding result.

The earlier the preceding theorem proves that
\(h_3\) is not monotone for
\[
0<a<\frac12
\qquad\text{or}\qquad
1<a<2.
\]
The middle-window theorem above proves increasing behavior for
\[
\frac12\le a\le1.
\]
Therefore Du--Wang Open Problem 2 is classified on \(0<a<2\):
\[
h_3 \text{ is increasing on }(0,\infty)
\quad\Longleftrightarrow\quad
\frac12\le a\le1,
\]
and it is not monotone on the two outer windows.

This proves the claim.
the corresponding result true by source-problem classification.
\end{proof}

The following appendix result records the fixed source problem \textbf{APP-0018} without interrupting the main exposition.
\begin{corollary}[APP-0018: Qi-Guo tau threshold has a supremum-corrected exact value]
\label{res:qg-tau-global-supremum}
Let \(a_*>0\) be the unique positive root of \(e^{a_*}=1+a_*+a_*^2\). Then \(\sup_{s\in\mathbb N,t>0}\tau(s,t)=a_*/(1+a_*+a_*^2)=0.298425607525639\ldots\), and the supremum is not attained.

This resolves the fixed application label \textbf{APP-0018}: Find the optimal uniform threshold for the tau expression in Open Problem 3.
\emph{Source reference.} \cite{app-app-0018-source}.
\end{corollary}
\begin{proof}
Writing \(u=t/(t+1)\) gives a two-parameter expression on \(0<u<1\).  The
Euler-Lagrange equations reduce the interior extremal sequence to the scaling
limit \(u=e^{-a/s}\).  The limiting profile is maximized exactly when
\(e^a=1+a+a^2\), and the resulting value is \(a/(1+a+a^2)\).  The source
one-variable monotonicity estimates bound the finite-\(s\) expressions by this
limit, while the scaling sequence approaches it.  Hence the number is a
supremum and is not attained at finite \(s,t\).
\end{proof}

The following appendix result records the fixed source problem \textbf{APP-0023} without interrupting the main exposition.
\begin{corollary}[APP-0023: Sokal generalized-Stieltjes lambda-derivatives have a triangular nonnegative compression]
\label{res:sokal-gs-lambda-derivative-generating-function}
For fixed \(n\), Sokal's generalized-Stieltjes tests satisfy the formal exponential generating function \(\sum_{k\ge0}F^{[\lambda]}_{n,k}(x)z^k/k!=(1-z)^{-(n+\lambda)}(-1)^n\sum_{j\ge0}x^j f^{(n+j)}(x)(z/(1-z))^j/j!\).

This resolves the fixed application label \textbf{APP-0023}: Explain the lambda-derivative structure of Sokal's generalized-Stieltjes Hankel-type expressions.
\emph{Source reference.} \cite{app-app-0023-source}.
\end{corollary}
\begin{proof}
The exponential generating function is
\[
 \mathcal F(z)
 =(1-z)^{-a}\sum_{j\ge0}G_j(x)
      \left({z\over 1-z}\right)^j{1\over j!}.
\]
Differentiating \(\ell\) times with respect to \(\lambda\) multiplies
\(\mathcal F\) by \((-\log(1-z))^\ell\).  Extracting coefficients gives the
displayed triangular formula.  Since
\(-\log(1-z)=\sum_{m\ge1}z^m/m\) has non-negative coefficients, all compression
coefficients are non-negative.
\end{proof}

\begin{proposition}
\label{res:gbf-product-closure}
If \(\lambda_1,\lambda_2>0\), \(f_1\in\mathcal B_{\lambda_1}\), and \(f_2\in\mathcal B_{\lambda_2}\), then \(f_1f_2\in\mathcal B_{\lambda_1+\lambda_2}\).
\end{proposition}
\begin{proof}
For \(\lambda>0\), the source defines \(f\in\mathcal B_\lambda\) when
\[
x^{1-\lambda}f'(x)
\]
is completely monotone. It also recalls generalized Stieltjes classes \(\mathcal S_\lambda\) and the product closure
\[
\mathcal S_{\lambda_1}\mathcal S_{\lambda_2}\subseteq \mathcal S_{\lambda_1+\lambda_2}.
\]

Proposition 5.1 states that if
\[
f_1\in\mathcal B_{\lambda_1},\qquad f_2\in\mathcal B_{\lambda_2},
\]
then
\[
f_1f_2\in\mathcal B_{\lambda_1+\lambda_2}.
\]

Local proof audit: by definition \(f_j'(x)/x^{\lambda_j-1}\) is completely monotone, and by the cited Corollary 2.1, \(f_j(x)/x^{\lambda_j}\) is completely monotone. Then
\[
\frac{(f_1f_2)'(x)}{x^{\lambda_1+\lambda_2-1}}
=
\frac{f_1'(x)}{x^{\lambda_1-1}}\frac{f_2(x)}{x^{\lambda_2}}
+
\frac{f_1(x)}{x^{\lambda_1}}\frac{f_2'(x)}{x^{\lambda_2-1}}.
\]
The right side is a positive sum of products of completely monotone functions, hence completely monotone.

The same source's Proposition 5.11 defines a class \(\mathcal C_\alpha\): \(f\in\mathcal C_\alpha\) when \(x^\alpha f(x)\) is the Laplace transform of a decreasing logarithmically convex function. It proves
\[
f\in\mathcal C_\alpha
\quad\Longrightarrow\quad
\frac1f\in\mathcal B_{\alpha+1}
\]
and
\[
\frac{1}{x^{\alpha+1}f(x)}
\]
is completely monotone.

Promote bridge nodes:

Demote/quarantine:
\end{proof}

The following appendix result records the fixed source problem \textbf{APP-0036} without interrupting the main exposition.
\begin{counterexample}[APP-0036: Baskakov even-line complete-monotonicity fails at \(r=8\)]
\label{res:not-baskakov-higher-power-even-conjecture}
Abel--Gawronski--Neuschel's even-power Baskakov complete-monotonicity conjecture is false; it fails at \(\alpha=1\), \(r=8\).

This resolves the fixed application label \textbf{APP-0036}: Decide the even-line Baskakov complete-monotonicity conjecture at \(r=8\).
\emph{Source reference.} \cite{app-app-0036-source}.
\end{counterexample}
\begin{proof}
Set
\[
D_m(x)=(1+x)^{2m}-x^{2m}.
\]
The roots are obtained from
\[
\frac{1+x}{x}=\omega_k,\qquad
\omega_k=e^{\pi i k/m},\qquad k=1,\ldots,2m-1.
\]
Thus
\[
\lambda_{k,m}
=\frac{1}{\omega_k-1}
=-\frac12-\frac{i}{2}\cot\frac{\pi k}{2m}.
\]
At this root,
\[
D_m'(\lambda_{k,m})
=2m\lambda_{k,m}^{2m-1}(\omega_k^{-1}-1),
\]
and direct simplification gives the real residue
\[
R_{k,m}
=\frac{1}{D_m'(\lambda_{k,m})}
=\frac{(-1)^{m+k}}{2m}
\left(2\sin\frac{\pi k}{2m}\right)^{2m-2}.
\]
Consequently the inverse-Laplace density is
\[
\rho_m(t)
=\sum_{k=1}^{2m-1}R_{k,m}e^{\lambda_{k,m}t}.
\]
Pairing conjugate roots gives
\[
\rho_m(t)
=\frac{e^{-t/2}}{2m}
\left[
2^{2m-2}
+2\sum_{k=1}^{m-1}
(-1)^{m+k}
\left(2\sin\frac{\pi k}{2m}\right)^{2m-2}
\cos\!\left(
\frac{t}{2}\cot\frac{\pi k}{2m}
\right)
\right].
\]

This proves the finite trigonometric density formula used by the diagnostic
candidate.

For \(m=2\), this gives
\[
\rho_2(t)=e^{-t/2}\left(1-\cos\frac t2\right)
=2e^{-t/2}\sin^2\frac t4\ge0,
\]
recovering the previously admitted \(r=4,\alpha=1\) seed.

For \(m=3\),
\[
\rho_3(t)
=\frac{e^{-t/2}}6
\left[
16+2\cos\left(\frac{\sqrt3\,t}{2}\right)
-18\cos\left(\frac{t}{2\sqrt3}\right)
\right].
\]
Set \(u=t/(2\sqrt3)\).  Since
\(\cos(3u)=4\cos^3u-3\cos u\),
\[
\rho_3(t)
=\frac{4}{3}e^{-t/2}(1-\cos u)^2(2+\cos u)\ge0.
\]
Thus \(r=6,\alpha=1\) is completely monotone by an explicit positive density.

For \(m=4\), write \(h_4(t)=e^{t/2}\rho_4(t)\).  The density formula gives
\[
h_4(t)
=8+2\cos\frac t2
-\frac{(2-\sqrt2)^3}{4}
\cos\left(\frac{1+\sqrt2}{2}t\right)
-\frac{(2+\sqrt2)^3}{4}
\cos\left(\frac{\sqrt2-1}{2}t\right).
\]
At \(t=10\pi\),
\[
\cos(t/2)=\cos(5\pi)=-1,
\]
and both other cosine terms equal \(-\cos(5\pi\sqrt2)\).  Since
\[
(2-\sqrt2)^3+(2+\sqrt2)^3=40,
\]
we obtain
\[
h_4(10\pi)=6+10\cos(5\pi\sqrt2).
\]
Now
\[
0<5\sqrt2-7<\frac14
\]
because \(7/5<\sqrt2<29/20\).  Hence, with
\(\delta=5\sqrt2-7\),
\[
\cos(5\pi\sqrt2)=\cos(7\pi+\pi\delta)=-\cos(\pi\delta)
<-\cos\frac\pi4=-\frac{\sqrt2}{2}<-\frac35.
\]
Therefore
\[
h_4(10\pi)<6+10\left(-\frac35\right)=0,
\]
and
\[
\rho_4(10\pi)=e^{-5\pi}h_4(10\pi)<0.
\]
Since the inverse-Laplace transform is unique for measures on
\([0,\infty)\), \(f^{[8]}_1(x)=1/((1+x)^8-x^8)\) cannot be completely monotone.
\end{proof}

The following appendix result records the fixed source problem \textbf{APP-0020} without interrupting the main exposition.
\begin{counterexample}[APP-0020: Qi's degree-four conjecture for h\_lambda is refuted at the endpoint]
\label{res:not-qi-hlambda-degree4-conjecture}
not(For \(\Psi(x)=[\psi'(x)]^2+\psi''(x)\) and \(h_\lambda(x)=\Psi(x)-(x^2+\lambda x+12)/(12x^4(x+1)^2)\), \(\deg_x^{\rm cm}h_\lambda=\deg_x^{\rm cm}(-h_\mu)=4\) if and only if \(\lambda\le0\) and \(\mu\ge4\).)

This resolves the fixed application label \textbf{APP-0020}: Test the conjectured complete-monotonic degree four for \(h\_\lambda\) and \(-h\_\mu\).
\emph{Source reference.} \cite{app-app-0020-source}.
\end{counterexample}
\begin{proof}
Near zero,
\[
 h_\lambda(x)=-{\lambda\over 12x^3}
 +\left({\lambda\over6}+{\pi^2\over3}-{37\over12}\right){1\over x^2}
 +O(x^{-1}).
\]
If \(\lambda<0\), then \(x^4h_\lambda(x)=(-\lambda/12)x+O(x^2)\), whose
derivative is positive near zero.  If \(\lambda=0\), the coefficient
\(\pi^2/3-37/12\) is positive, again giving a positive derivative for the
degree-four transform near zero.  Complete monotonicity would require the first
derivative to be non-positive.  The same small-\(x\) test applied to
\(-x^4h_\mu\) excludes the claimed \(\mu\)-range.  Hence the degree-four
conjecture fails.
\end{proof}

The following appendix result records the fixed source problem \textbf{APP-0013} without interrupting the main exposition.
\begin{corollary}[APP-0013: Yin \((p,k)\)-digamma sharp \(\alpha\)-necessity]
\label{res:yin-pk-digamma-alpha-necessity}
For \(p\in\mathbb N\), \(k>0\), if \(x^\alpha\left[\frac1k\log\frac{pkx}{x+k(p+1)}-\psi_{p,k}(x)\right]\) is completely monotone on \((0,\infty)\), then \(\alpha\le1\).

This resolves the fixed application label \textbf{APP-0013}: Yin \((p,k)\)-digamma sharp \(\alpha\)-necessity
\emph{Source reference.} \cite{app-app-0013-source}.
\end{corollary}
\begin{proof}
\emph{Setup.}
Fix \(p\in\mathbb N\) and \(k>0\). Define
\[
B_{p,k}(x)=
\frac1k\log\frac{pkx}{x+k(p+1)}-\psi_{p,k}(x),
\qquad
\delta_{p,k,\alpha}(x)=x^\alpha B_{p,k}(x).
\]

The source gives the finite formula
\[
\psi_{p,k}(x)=\frac1k\log(pk)-\sum_{n=0}^p\frac1{x+nk}.
\]
Therefore
\[
\begin{aligned}
B_{p,k}(x)
&=
\frac1k\log\frac{pkx}{x+k(p+1)}
-\frac1k\log(pk)
+\sum_{n=0}^p\frac1{x+nk}  \\
&=
\frac1k\log\frac{x}{x+k(p+1)}
+\sum_{n=0}^p\frac1{x+nk}.
\end{aligned}
\]

This normal form also gives positivity directly. Put \(t=x/k\). Then
\[
B_{p,k}(x)
=\frac1k\left[
\sum_{n=0}^p\frac1{t+n}
+\log\frac{t}{t+p+1}
\right].
\]
Since
\[
\log\frac{t+p+1}{t}=\int_0^{p+1}\frac{du}{t+u},
\]
and \(u\mapsto(t+u)^{-1}\) is strictly decreasing, the left Riemann sum dominates:
\[
\sum_{n=0}^p\frac1{t+n}
>
\int_0^{p+1}\frac{du}{t+u}.
\]
Hence \(B_{p,k}(x)>0\) for \(x>0\).

This proves the claim.

From the finite normal form,
\[
B_{p,k}(x)
=\frac1x
+\frac1k\log\frac{x}{x+k(p+1)}
+\sum_{n=1}^p\frac1{x+nk}.
\]

As \(x\to0^+\), the finite sum over \(n\ge1\) is \(O(1)\), and
\[
\frac1k\log\frac{x}{x+k(p+1)}=O(|\log x|).
\]
Thus
\[
B_{p,k}(x)=\frac1x+O(|\log x|)
\qquad (x\to0^+).
\]
Equivalently,
\[
xB_{p,k}(x)\to1.
\]
Therefore
\[
\delta_{p,k,\alpha}(x)
=x^\alpha B_{p,k}(x)
=x^{\alpha-1}(xB_{p,k}(x))
\sim x^{\alpha-1}.
\]

This proves the claim.

Assume that \(\delta_{p,k,\alpha}\) is completely monotone on \((0,\infty)\). Then
\[
\delta_{p,k,\alpha}(x)\ge0,\qquad
\delta_{p,k,\alpha}'(x)\le0.
\]
Since \(B_{p,k}(x)>0\), the function \(\delta_{p,k,\alpha}\) is actually positive for every \(x>0\).

Suppose, for contradiction, that \(\alpha>1\). The endpoint asymptotic gives
\[
\lim_{x\to0^+}\delta_{p,k,\alpha}(x)=0.
\]
Fix \(x_0>0\). Positivity gives \(\delta_{p,k,\alpha}(x_0)>0\). Because \(\delta_{p,k,\alpha}(x)\to0\) as \(x\to0^+\), choose \(0<y<x_0\) with
\[
\delta_{p,k,\alpha}(y)<\frac12\delta_{p,k,\alpha}(x_0).
\]
But \(\delta_{p,k,\alpha}'\le0\) means the function is nonincreasing as \(x\) increases, hence for \(0<y<x_0\),
\[
\delta_{p,k,\alpha}(y)\ge\delta_{p,k,\alpha}(x_0),
\]
a contradiction.

Therefore \(\alpha\le1\).

This proves the claim.
\end{proof}

The following appendix result records the fixed source problem \textbf{APP-0039} without interrupting the main exposition.
\begin{counterexample}[APP-0039: Ramanujan Turan-window complete-monotonicity interval is false]
\label{res:ramanujan-turan-window-negative-answer}
The Mishra--Swaminathan Ramanujan integral Turan-window complete-monotonicity problem has a negative answer: the proposed interval already fails for \(n=2\) and suitable \(\alpha\in(0,1/2)\).

This resolves the fixed application label \textbf{APP-0039}: Decide the Ramanujan Turan-window complete-monotonicity interval.
\emph{Source reference.} \cite{app-app-0039-source}.
\end{counterexample}
\begin{proof}
Put
\[
A_k(x)=(-1)^k I_R^{(k)}(x)
=\int_0^\infty e^{-xt}\frac{t^{k-1}}{\pi^2+\log^2 t}\,dt
\qquad (k\ge1).
\]
Then
\[
H_2(x;\alpha)=A_2(x)^2-\alpha A_1(x)A_3(x).
\]

We study \(x=e^{-L}\) as \(L\to\infty\). With \(y=xt\),
\[
A_k(e^{-L})
=
e^{kL}\int_0^\infty
\frac{e^{-y}y^{k-1}}
{\pi^2+(L+\log y)^2}\,dy.
\]

For fixed \(k=1,2,3\),
\[
A_k(e^{-L})
=
e^{kL}L^{-2}\Gamma(k)
\left(1-\frac{2\psi(k)}{L}+O(L^{-2})\right),
\]
where \(\psi\) is the digamma function. To justify the expansion, split the \(y\)-integral into the central region \(|\log y|\le L^{2/3}\) and its complement. On the central interval, with \(z=\log y\), expand uniformly:
\[
\frac{1}{\pi^2+(L+\log y)^2}
=
L^{-2}\left(1-\frac{2\log y}{L}+O\left(\frac{1+\log^2 y}{L^2}\right)\right),
\]
and integrate against \(e^{-y}y^{k-1}\). The lower complement \(0<y<e^{-L^{2/3}}\) is \(O(e^{-cL^{2/3}})\) for \(k=1,2,3\) after the change \(y=e^{-u}\), and the upper complement \(y>e^{L^{2/3}}\) is super-exponentially small because of \(e^{-y}\). These errors are \(o(L^{-N})\) for every fixed \(N\), which is more than needed for the displayed expansion.

Therefore
\[
\frac{A_2(e^{-L})^2}{A_1(e^{-L})A_3(e^{-L})}
=
\frac{\Gamma(2)^2}{\Gamma(1)\Gamma(3)}
\left(
1+\frac{2(\psi(1)+\psi(3)-2\psi(2))}{L}+O(L^{-2})
\right).
\]
Using
\[
\psi(1)=-\gamma,\qquad
\psi(2)=1-\gamma,\qquad
\psi(3)=\frac32-\gamma,
\]
we get
\[
\psi(1)+\psi(3)-2\psi(2)=-\frac12,
\]
so
\[
\frac{A_2(e^{-L})^2}{A_1(e^{-L})A_3(e^{-L})}
=
\frac12-\frac{1}{2L}+O(L^{-2}).
\]

Choose
\[
\alpha_L=\frac12-\frac{1}{4L}.
\]
For all sufficiently large \(L\), \(\alpha_L\in(0,1/2)\) and
\[
\frac{A_2(e^{-L})^2}{A_1(e^{-L})A_3(e^{-L})}<\alpha_L.
\]
Consequently
\[
H_2(e^{-L};\alpha_L)
=
A_2(e^{-L})^2-\alpha_L A_1(e^{-L})A_3(e^{-L})
<0.
\]

Complete monotonicity would imply nonnegativity at order zero. Hence \(H_2(\cdot;\alpha_L)\) is not completely monotone for these fixed parameters \(\alpha_L\in(0,1/2)\). The source's universal Turan-window statement is false already in the \(n=2\) upper window.

The new mechanism is a logarithmic-density moment-ratio asymptotic obstruction for quadratic Laplace-moment Turan gaps. It is distinct from the public APP-0014 Stieltjes/complete-Bernstein classification of \(I_R\).
\end{proof}

The following appendix result records the fixed source problem \textbf{APP-0031} without interrupting the main exposition.
\begin{counterexample}[APP-0031: Baricz gamma-quotient Bernstein test is false]
\label{res:baricz-gamma-quotient-bf-forall-negative-answer}
Baricz's question asking whether \(x\mapsto\Gamma(x)\Gamma(x-a+b)/(\Gamma(x-a)\Gamma(x+b))\) is a Bernstein function on \((a,\infty)\) for every \(a,b>0\) has a negative answer.

This resolves the fixed application label \textbf{APP-0031}: Decide whether Baricz's gamma-quotient is Bernstein for every \(a,b>0\).
\emph{Source reference.} \cite{app-app-0031-source}.
\end{counterexample}
\begin{proof}
Take \(a=2\), \(b=3\), and set \(y=x-2>0\). Then by the Gamma recurrence,
\[
\frac{\Gamma(x)\Gamma(x-a+b)}{\Gamma(x-a)\Gamma(x+b)}
=\frac{\Gamma(x)\Gamma(x+1)}{\Gamma(x-2)\Gamma(x+3)}
=\frac{(x-2)(x-1)}{(x+1)(x+2)}
=\frac{y(y+1)}{(y+3)(y+4)}.
\]

Let
\[
g(y)=\frac{y(y+1)}{(y+3)(y+4)},\qquad y>0.
\]

If \(g\) were a Bernstein function, then \(g'\) would be completely monotone. In particular \(g'''\ge0\) would be necessary.

Exact symbolic differentiation gives
\[
g'(y)=\frac{6(y^2+4y+2)}{(y+3)^2(y+4)^2},
\]
\[
g''(y)=
-\frac{12(y^3+6y^2+6y-10)}{(y+3)^3(y+4)^3},
\]
and
\[
g'''(y)=
\frac{36(y^4+8y^3+12y^2-40y-94)}{(y+3)^4(y+4)^4}.
\]

At \(y=1\),
\[
g'''(1)=-\frac{1017}{40000}<0.
\]

Therefore \(g'\) is not completely monotone, \(g\) is not Bernstein, and the Baricz for-all-\(a,b\) Bernstein question has a negative answer.
\end{proof}

\begin{proposition}
\label{res:bz-gamma-quotient-critical-window-reduction}
There is a rational interval \(J\) around \(1.126207061\ldots\) such that \(\widetilde F'<0\) on \((-1,\inf J]\), \(\widetilde F'>0\) on \([\sup J,\infty)\), and \(\widetilde F'\) has exactly one zero inside \(J\).
\end{proposition}
\begin{proof}
\emph{Setup.}
Use the notation from the previous Bulboaca--Zayed pass:
\[
G(x)=\log\Gamma(x+1),\qquad
D(x)=\log\frac{x^2+6}{x+6},
\]
and
\[
F(x)=\frac{G(x)}{D(x)}.
\]
For \(x>1\), \(D(x)>0\), and
\[
D'(x)=p(x)=\frac{x^2+12x-6}{(x^2+6)(x+6)}.
\]
Let
\[
N_0(x)=\psi(x+1)D(x)-G(x)D'(x).
\]
The derivative-sign numerator recorded in
the corresponding result is
\[
N(x)=((x^2+6)(x+6))\,N_0(x),
\]
so \(N\) and \(N_0\) have the same sign on \(x>1\).

The source itself proves the decreasing part on \((-1,1)\), and the earlier
the preceding theorem proves \(F'(x)>0\) for
\(x\ge8\).  It remains to prove a single sign change of \(N_0\) on \((1,8]\).

Put
\[
P(x)=x^2+12x-6,\qquad Q(x)=(x^2+6)(x+6),
\]
so \(p=P/Q\), and define
\[
r(x)=\frac{\psi(x+1)}{p(x)}=\psi(x+1)\frac{Q(x)}{P(x)}.
\]
Then
\[
P(x)^2r'(x)=S(x),
\]
where
\[
S(x)=\psi'(x+1)Q(x)P(x)+\psi(x+1)K(x),
\]
and
\[
K(x)=Q'(x)P(x)-Q(x)P'(x)
=x^4+24x^3+48x^2-144x-468.
\]
Direct differentiation gives
\[
S'(x)=P(x)\{Q(x)\psi''(x+1)+2Q'(x)\psi'(x+1)\}+K'(x)\psi(x+1).
\]

For \(y=x+1\ge2\), use the standard elementary bounds
\[
\psi(y)>\log y-\frac1y,\qquad
\psi'(y)>\frac1y+\frac{1}{2y^2},\qquad
\psi''(y)>-\frac1{y^2}-\frac2{y^3}.
\]
The last inequality follows from the series
\[
\psi''(y)=-2\sum_{n=0}^{\infty}(n+y)^{-3}
\]
and an integral upper bound for the sum.

Substituting these bounds into \(S'\) yields
\[
S'(x)>
P(x)\frac{5x^4+30x^3+57x^2+12x-90}{(x+1)^3}
K'(x)\left(\log(x+1)-\frac1{x+1}\right).
\]
On \(x\ge1\), \(P(x)>0\).  The quartic numerator has value \(14\) at \(x=1\)
and positive derivative
\[
20x^3+90x^2+114x+12.
\]
Also
\[
K'(x)=4x^3+72x^2+96x-144
\]
has value \(28\) at \(x=1\) and positive derivative for \(x\ge1\), while
\(\log(x+1)-1/(x+1)>0\) for \(x\ge1\).  Hence
\[
S'(x)>0\qquad(1\le x\le8).
\]

At the endpoints,
\[
S(1)=343\left(\frac{\pi^2}{6}-1\right)-539(1-\gamma)<0,
\]
for example using \(\pi^2<987/100\) and \(\gamma<5773/10000\), which gives
\(-66003/10000<0\).  Also \(S(8)>0\), because \(K(8)=17836>0\),
\(\psi'(9)>0\), and \(\psi(9)=H_8-\gamma>0\).  Therefore \(S\) has exactly
one zero on \([1,8]\).  Since \(P^2>0\), \(r'\) has exactly one zero; \(r\)
decreases first and then increases.

Now define
\[
I(x)=D(x)r(x)-G(x).
\]
Then
\[
N_0(x)=p(x)I(x),\qquad I'(x)=D(x)r'(x).
\]
Because \(D(x)>0\) for \(x>1\), \(I\) decreases and then increases with the
same turning point as \(r\).  Moreover \(\lim_{x\downarrow1}I(x)=0\), so
\(I(x)<0\) immediately to the right of \(1\).

Finally,
\[
N_0(8)=(H_8-\gamma)\log5-\frac{11}{70}\log(40320)>0.
\]
For instance, the bounds
\[
\gamma<\frac{5773}{10000},\quad \log5>\frac{1609}{1000},
\quad \log(40320)<\frac{10605}{1000}
\]
give the lower bound
\[
\left(\frac{761}{280}-\frac{5773}{10000}\right)\frac{1609}{1000}
-\frac{11}{70}\frac{10605}{1000}
=\frac{124435951}{70000000}>0.
\]
Thus \(I\), and therefore \(N_0\) and \(N\), has exactly one zero on
\((1,8]\).  Denote it by \(\xi\).  The sign pattern is
\[
N(x)<0\quad(1<x<\xi),\qquad
N(\xi)=0,\qquad
N(x)>0\quad(\xi<x\le8).
\]

A high-precision numerical check places
\[
\xi=1.12620706105164346558\ldots,
\]
matching the source's reported value.  The numerical value is not used in the
proof.

The raw numerator has a removable zero at \(x=1\), since \(G(1)=D(1)=0\).  Its
quadratic leading coefficient is
\[
\lim_{x\downarrow1}\frac{N_0(x)}{(x-1)^2}
=
\frac{7(\pi^2/6-1)-11(1-\gamma)}{98}<0,
\]
again by \(\pi^2<987/100\) and \(\gamma<5773/10000\), which gives
\(-1347/980000<0\).  Therefore any normalization dividing out the removable
\((x-1)^2\) factor is negative at the left endpoint.

The single-crossing candidate is true:
\[
T\text{-BZ-gamma-quotient-N-single-crossing-one-eight}.
\]
Together with the source's theorem on \((-1,1)\) and the previous local
right-tail theorem for \(x\ge8\), this proves the critical-window reduction
\[
T\text{-BZ-gamma-quotient-critical-window-reduction}
\]
and then the source problem
\[
T\text{-Bulboaca-Zayed-gamma-quotient-full-monotonicity}.
\]

The finite-envelope candidate is not separately promoted: the proof above is a
clean analytic ratio-kernel argument rather than the finite rational cover
specified in that candidate.  The diagnostic obstruction-map candidate is also
not needed after the single-crossing proof.
\end{proof}

The following appendix result records the fixed source problem \textbf{APP-0011} without interrupting the main exposition.
\begin{corollary}[APP-0011: Exact \(n=2\) beta-window endpoint]
\label{res:q2-terminal-exact}
The exact admissible beta set \(\mathcal I_2\) for Qi--Lim--Nantomah Open Problem 4 is determined by a certified value or certified description of \(L_2=\sup_{0<x<1}Q_2(x)\), including the endpoint-inclusion convention.

This resolves the fixed application label \textbf{APP-0011}: Exact \(n=2\) beta-window endpoint
\emph{Source reference.} \cite{app-app-0011-source}.
\end{corollary}
\begin{proof}
This roll proves a true auxiliary finite-middle outside-cover node:

the corresponding result.

The replayable certificate is:

\begin{verbatim}
\end{verbatim}

It proves
\[
Q_2(x)<q_J
\]
on
\[
\left[\frac9{50},\frac{1409}{5000}\right]
\cup
\left[\frac{293}{1000},\frac9{10}\right].
\]

The left bridge \([9/50,221050/1000000]\) uses the already-derived identity
\[
3xZ_4(x)-Z_3(x)
=2x^{-3}+\sum_{k=1}^{\infty}\frac{2x-k}{(x+k)^4}.
\]
For \(x<1/2\), this gives
\[
R(x)>Z_3(1/x).
\]
The certificate proves
\[
Z_3(1/x)>x^{1157/500}
\]
on \([9/50,221050/1000000]\). Since
\[
\frac{1157}{500}<\frac{115727}{50000}<q_J,
\]
this gives \(Q_2(x)<q_J\) on the bridge.

The remaining two subintervals use the direct comparison
\[
\log R(x)>\frac{115727}{50000}\log x.
\]
Since \(\log x<0\), this implies
\[
Q_2(x)<\frac{115727}{50000}<q_J.
\]
The script uses power-of-two range reduction for the atanh logarithm enclosure:
\[
\log r=\log(2^k r)-k\log 2,
\]
so every logarithm is evaluated with an argument in \([1,2)\).

The certificate output was:

\begin{verbatim}
REMAINING_FINITE_MIDDLE_COVER_OK
ZETA_LEFT_SUBCOVER
interval_count 245
domain [180000/1000000,221050/1000000]
min_width 1 max_width 1283
worst_interval [180000/1000000,181282/1000000]
worst_margin_positive True
DIRECT_LEFT_SUBCOVER
interval_count 3739
domain [221050/1000000,281800/1000000]
min_width 1 max_width 238
worst_interval [280389/1000000,280392/1000000]
worst_margin_positive True
DIRECT_RIGHT_SUBCOVER
interval_count 3801
domain [293000/1000000,900000/1000000]
min_width 1 max_width 18969
worst_interval [293124/1000000,293126/1000000]
worst_margin_positive True
\end{verbatim}

Together with the earlier true covers \((0,9/50]\) and \([9/10,1)\), this covers all outside-compact points below \(q_J\).

This roll also proves:

the corresponding result.

The replayable certificate is:

\begin{verbatim}
\end{verbatim}

Let
\[
H(x)=\Lambda(x)+x\Lambda'(x).
\]
Using \(Z_s'(x)=-sZ_{s+1}(x)\) and
\[
\frac{d}{dx}Z_s(1/x)=s x^{-2} Z_{s+1}(1/x),
\]
the script builds a rational interval enclosure for \(H\) using \(Z_3,\ldots,Z_6\) at \(x\) and \(Z_3,\ldots,Z_5\) at \(1/x\). It proves
\[
H(x)>0
\qquad
\left(\frac{1409}{5000}\le x\le\frac{293}{1000}\right).
\]

Since
\[
G'(x)=\log x\,H(x)
\]
and \(\log x<0\) on the compact bracket, this gives
\[
G'(x)<0.
\]
The existing endpoint certificate proves
\[
G(287345/1000000)>0,
\qquad
G(287346/1000000)<0.
\]
Therefore \(G\) has a unique zero
\[
\xi\in J=\left[\frac{287345}{1000000},\frac{287346}{1000000}\right],
\]
with \(G>0\) to the left of \(J\) and \(G<0\) to the right of \(J\) on the compact bracket.

Because
\[
Q_2'(x)=\frac{G(x)}{x(\log x)^2},
\]
\(Q_2\) is increasing before \(\xi\) and decreasing after \(\xi\) on the compact bracket.

The certificate output was:

\begin{verbatim}
COMPACT_MONOTONICITY_CERTIFICATE_OK
endpoint_signs G_left_positive G_right_negative
h_interval_count 128
domain [281800/1000000,293000/1000000]
min_width 87 max_width 88
worst_interval [281887/1000000,281975/1000000]
worst_H_lower_positive True
\end{verbatim}

The true cover package is now:

\[
(0,9/50]\cup
\left[9/50,\frac{1409}{5000}\right]\cup
\left[\frac{293}{1000},9/10\right]\cup
[9/10,1).
\]

All points outside the compact bracket have \(Q_2(x)<q_J\), and \(q_J<Q_2(\xi)\) by the certified inner witness. On the compact bracket, \(Q_2\) has a unique maximum at the unique zero \(\xi\in J\). Thus
\[
L_2=Q_2(\xi).
\]

At \(\beta=Q_2(\xi)\), the defining inequality has equality at \(x=\xi\), so the lower endpoint is excluded. The known upper endpoint remains included. Therefore
\[
\mathcal I_2=(Q_2(\xi),3].
\]

This proves the terminal exact endpoint node.
\end{proof}

The following appendix result records the fixed source problem \textbf{APP-0019} without interrupting the main exposition.
\begin{corollary}[APP-0019: Wakrim W-symbol Bernstein range closes the \(\beta\)-gap]
\label{res:wakrim-w-operator-bf-gap-solved}
Wakrim's open Bernstein-status gap for the W-symbol in the range \(0<\beta\le1\) is resolved affirmatively: \(\Phi_{\alpha,\beta}\) is Bernstein for every \(0<\alpha<1\) and \(0<\beta\le1\).

This resolves the fixed application label \textbf{APP-0019}: Close the open Bernstein-symbol range for the W-operator when \(0<\beta\le1\).
\emph{Source reference.} \cite{app-app-0019-source}.
Wakrim's source supplies the non-Bernstein obstruction for \(\beta>1\); the present proof supplies the remaining Bernstein range \(0\le\beta\le1\). No complete-Bernstein or subordination conclusion is asserted here.
\end{corollary}
\begin{proof}
Wakrim's source supplies the non-Bernstein obstruction for \(\beta>1\); the present proof supplies the remaining Bernstein range \(0\le\beta\le1\). No complete-Bernstein or subordination conclusion is asserted here.

Wakrim's source proves the non-Bernstein obstruction for \(\beta>1\) and
leaves the range \(0<\beta\le1\) as the remaining delicate case.  We prove
that remaining Bernstein half in detail.  No complete-Bernstein or
subordination conclusion is asserted here.

Put
\[
  a=1-\alpha\in(0,1),\qquad y=s^a .
\]
Since \(\alpha=1-a\),
\[
 \Phi_{\alpha,\beta}(s)
 =s^{1-a}(1+a s^{-a})^{-\beta}
 =s^{1-a+a\beta}(s^a+a)^{-\beta}.
\]
Let
\[
  p=1-a+a\beta=\alpha+a\beta .
\]
Differentiating the last expression gives
\[
\begin{aligned}
 \Phi'_{\alpha,\beta}(s)
 &=s^{p-1}(s^a+a)^{-\beta-1}
     \{p(s^a+a)-\beta a s^a\}  \\
 &=s^{p-1}(s^a+a)^{-\beta-1}
     \{\alpha s^a+\alpha a+a^2\beta\}       \\
 &=s^{p-1}(s^a+a)^{-\beta}
     \left(\alpha+{a^2\beta\over s^a+a}\right).
\end{aligned}
\]
Because \(p-1=a(\beta-1)\), this is
\[
  \Phi'_{\alpha,\beta}(s)=H_{\alpha,\beta}(s^a),
\]
where
\[
 H_{\alpha,\beta}(y)
 =y^{\beta-1}(y+a)^{-\beta}
      \left(\alpha+{a^2\beta\over y+a}\right).
\]

We now check complete monotonicity of \(H_{\alpha,\beta}\) for
\(0\le\beta\le1\).  The first factor is
\[
  y^{\beta-1}=y^{-(1-\beta)},
\]
and this is interpreted as follows.  If \(0\le\beta<1\), put
\(r=1-\beta>0\).  Then
\[
  y^{\beta-1}=y^{-r}
  ={1\over\Gamma(r)}
    \int_0^\infty e^{-yt}t^{r-1}\,dt,
\]
so \(y^{\beta-1}\) is completely monotone by the Bernstein--Widder
criterion.  If \(\beta=1\), the same factor is the constant \(1\).  The shifted factor
\((y+a)^{-\beta}\) is completely monotone for \(\beta>0\) by the Laplace
representation
\[
 (y+a)^{-\beta}
 ={1\over\Gamma(\beta)}
   \int_0^\infty e^{-yt}e^{-at}t^{\beta-1}\,dt,
\]
and for \(\beta=0\) it is the constant \(1\).  The last factor
\[
  \alpha+{a^2\beta\over y+a}
\]
is a positive constant plus a nonnegative multiple of the completely monotone
resolvent \((y+a)^{-1}\).  By closure of completely monotone functions under
products and nonnegative sums, \(H_{\alpha,\beta}\) is completely monotone.

The endpoint \(\beta=1\) has no limiting ambiguity; explicitly
\[
 H_{\alpha,1}(y)
 ={\alpha\over y+a}+{a^2\over (y+a)^2},
\]
which is completely monotone.  The endpoint \(\beta=0\) also agrees with the
direct identity \(\Phi_{\alpha,0}(s)=s^\alpha\).

Finally, let \(g(s)=s^a\).  The composition theorem used here is the standard
one: if \(f\) is completely monotone on \((0,\infty)\) and \(g\) is a
Bernstein function with \(g((0,\infty))\subset(0,\infty)\), then \(f\circ g\)
is completely monotone.  Its hypotheses are satisfied in the present case:
\(H_{\alpha,\beta}\) is completely monotone by the factor check above,
\(g(s)=s^a>0\), and \(g\) is a Bernstein function for \(0<a<1\) since
\[
  g'(s)=a s^{a-1}=a s^{-(1-a)}
\]
is completely monotone.  Hence
\(\Phi'_{\alpha,\beta}=H_{\alpha,\beta}\circ g\) is completely monotone for
\(0\le\beta\le1\), so \(\Phi_{\alpha,\beta}\) is a Bernstein function.
Together with Wakrim's source obstruction for \(\beta>1\), this closes the
remaining case and gives the exact Bernstein range.
\end{proof}

The following appendix result records the fixed source problem \textbf{APP-0006} without interrupting the main exposition.
\begin{corollary}[APP-0006: Exact inverse-tail floor formula at s=8]
\label{res:zeta-tail-floor-next-case}
Find a new exact reciprocal zeta-tail floor formula beyond the staged \(s=7\) and \(s=8\) cases, or prove a reusable enclosure template that yields such formulas for further integer exponents.

This resolves the fixed application label \textbf{APP-0006}: Exact inverse-tail floor formula at \(s=8\)
\emph{Source reference.} \cite{app-app-0006-source}.
\end{corollary}
\begin{proof}
For integer \(s>1\), the Hurwitz-tail Euler--Maclaurin expansion has the form
\[
\zeta_n(s)
=n^{1-s}\left[
\frac1{s-1}+\frac{z}{2}
+\sum_{r\ge1}\frac{B_{2r}}{(2r)!}(s)_{2r-1}z^{2r}
\right],
\qquad z=\frac1n.
\]
Formally inverting the bracket gives
\[
\zeta_n(s)^{-1}\sim n^{s-1}G_s(z).
\]

The replay script
computes \(G_s\) symbolically and verifies two things:

1. For \(s=8\), the coefficients through \(z^{15}\) reproduce the staged \(A_8\) and \(B_8-A_8\) coefficients.
2. For \(s=9\), truncating \(n^8G_9(1/n)\) through \(z^{14}\) gives an \(A_9\), and adding the \(z^{15}\) term gives \(B_9\), with exact telescoping residuals of the expected signs.

The generated \(s=9\) lower approximant is
\[
\begin{aligned}
A_9(n)=&8n^8-32n^7+80n^6-128n^5+120n^4-64n^3
+\frac{624}{7}n^2-\frac{512}{7}n-\frac{3324}{7}\\
&+\frac{90304}{21n^2}+\frac{90304}{21n^3}
-\frac{10280944}{245n^4}-\frac{64846304}{735n^5}
+\frac{109940976}{245n^6}.
\end{aligned}
\]
The upper approximant is
\[
B_9(n)=A_9(n)+\frac{384388288}{245n^7}.
\]

Let \(R_A(n)=1/A_9(n)\) and \(R_B(n)=1/B_9(n)\). Exact denominator clearing gives:
\[
R_A(k)-R_A(k+1)-\frac1{k^9}>0
\qquad(k\ge5),
\]
because the numerator, after substituting \(k=m+5\), has all coefficients positive. Similarly,
\[
\frac1{k^9}-\left(R_B(k)-R_B(k+1)\right)>0
\qquad(k\ge1),
\]
because the denominator-cleared numerator, after substituting \(k=m+1\), has all coefficients positive.

Therefore, for \(n\ge5\),
\[
A_9(n)<\zeta_n(9)^{-1}<B_9(n).
\]

This proves the inverse-window part of the reusable template and promotes the corresponding result.

The replay script also verifies the no-integer gap.

Write \(P_9(n)\) for the polynomial part of \(A_9(n)\):
\[
P_9(n)=8n^8-32n^7+80n^6-128n^5+120n^4-64n^3
+\frac{624}{7}n^2-\frac{512}{7}n-\frac{3324}{7}.
\]
Modulo \(7\), the fractional numerator of \(P_9(n)\) is
\[
n^2-n+1\pmod 7,
\]
so the fractional part of \(P_9(n)\) is one of
\[
0,\quad \frac17,\quad \frac37,\quad \frac67.
\]

The correction \(A_9(n)-P_9(n)\) is positive for \(n\ge9\); after denominator clearing and substituting \(n=m+9\), the numerator has all coefficients positive. For \(n\ge174\), the upper correction is bounded by
\[
B_9(n)-P_9(n)
<
\frac{90304}{21n^2}
+\frac{90304}{21n^3}
+\frac{109940976}{245n^6}
+\frac{384388288}{245n^7}
<\frac17.
\]
Therefore \(A_9(n)\) and \(B_9(n)\) cannot straddle an integer for \(n\ge174\). Exact rational arithmetic verifies
\[
\lfloor A_9(n)\rfloor=\lfloor B_9(n)\rfloor
\qquad(9\le n\le173).
\]
Consequently,
\[
\left\lfloor\zeta_n(9)^{-1}\right\rfloor=\lfloor A_9(n)\rfloor
\qquad(n\ge9).
\]

For \(1\le n\le8\), the exact values are
\[
\begin{array}{c|rrrrrrrr}
n&1&2&3&4&5&6&7&8\\
\hline
\lfloor\zeta_n(9)^{-1}\rfloor&0&497&18093&224086&1543530&7360226&27295287&84349541.
\end{array}
\]
For \(n=1\), this follows from \(\zeta(9)>1\). For \(2\le n\le8\), the replay script uses exact rational partial sums through \(N=100\) and the integral tail bound
\[
\sum_{k=N+1}^{\infty}k^{-9}<\frac{1}{8N^8}
\]
to prove
\[
\frac1{M_n+1}<\zeta_n(9)<\frac1{M_n}.
\]

This proves the claim.
\end{proof}

\begin{proposition}
\label{res:baskakov-diagonal-r4-c-equals-n-cm}
For every positive integer \(n\), the diagonal generalized Baskakov sum \(\psi^{[4]}_{n,n}(x)=\sum_{k=0}^{\infty}(p^{[n]}_{n,k}(x))^4\) is completely monotone on \((0,\infty)\).
\end{proposition}
\begin{proof}
For \(\alpha=1\),
\[
\binom{-1}{k}=(-1)^k,
\qquad
\binom{-1}{k}^{4}=1.
\]
Therefore, for \(x>0\),
\[
f^{[4]}_1(x)
=(1+x)^{-4}\sum_{k=0}^{\infty}
\left(\frac{x}{1+x}\right)^{4k}
=\frac{1}{(1+x)^4-x^4}.
\]
Since
\[
(1+x)^4-x^4=(2x+1)(2x^2+2x+1),
\]
partial fractions give
\[
\frac{1}{(1+x)^4-x^4}
=\frac{2}{2x+1}-\frac{2x+1}{2x^2+2x+1}.
\]
Writing \(a=x+\frac12\), this is
\[
\frac{1}{a}-\frac{a}{a^2+(1/2)^2}.
\]
Using
\[
\int_0^\infty e^{-xt}e^{-t/2}\,dt=\frac{1}{x+1/2}
\]
and
\[
\int_0^\infty e^{-xt}e^{-t/2}\cos(t/2)\,dt
=\frac{x+1/2}{(x+1/2)^2+(1/2)^2},
\]
we obtain
\[
f^{[4]}_1(x)
=\int_0^\infty e^{-xt}e^{-t/2}\left(1-\cos\frac{t}{2}\right)\,dt.
\]
The density \(e^{-t/2}(1-\cos(t/2))\) is nonnegative on \([0,\infty)\). By Bernstein's theorem, \(f^{[4]}_1\) is completely monotone on \((0,\infty)\).

From the source definition
\[
p^{[c]}_{n,k}(x)=\binom{-n/c}{k}(-cx)^k(1+cx)^{-n/c-k},
\]
one has
\[
\psi^{[r]}_{n,c}(x)=\sum_{k=0}^{\infty}\left(p^{[c]}_{n,k}(x)\right)^r
=f^{[r]}_{n/c}(cx)
\]
for even \(r\). Hence, for \(c=n\),
\[
\psi^{[4]}_{n,n}(x)=f^{[4]}_1(nx).
\]
Positive scaling preserves complete monotonicity, so \(\psi^{[4]}_{n,n}\) is completely monotone for every positive integer \(n\).
\end{proof}

The following appendix result records the fixed source problem \textbf{APP-0033} without interrupting the main exposition.
\begin{corollary}[APP-0033: Ramanujan antiderivative is complete Bernstein]
\label{res:ramanujan-antiderivative-complete-bernstein}
The antiderivative \(\widetilde I_R(x)=a+\int_0^\infty(1-e^{-xt})\,dt/[t(\pi^2+\log^2 t)]\) of the Ramanujan integral is a complete Bernstein function.

This resolves the fixed application label \textbf{APP-0033}: Resolve whether the Ramanujan integral antiderivative is a complete Bernstein function.
\emph{Source reference.} \cite{app-app-0033-source}.
\end{corollary}
\begin{proof}
The density
\[
\phi_0(t)=\frac{1}{t(\pi^2+\log^2 t)}
\]
is completely monotone on \((0,\infty)\). Consequently \(I_R\) is a Stieltjes function. Moreover, the antiderivative
\[
\widetilde I_R(x)=a+\int_0^\infty (1-e^{-xt})\phi_0(t)\,dt
\]
is a complete Bernstein function.

For \(u=\log t\),
\[
\int_0^1 e^{-au}\sin(\pi a)\,da
=\frac{\pi(1+e^{-u})}{u^2+\pi^2}.
\]
Substituting \(u=\log t\) gives
\[
\int_0^1 t^{-a}\sin(\pi a)\,da
=\frac{\pi(1+t^{-1})}{\pi^2+\log^2 t}.
\]
Therefore
\[
\phi_0(t)
=\frac{1}{t(\pi^2+\log^2 t)}
=\frac{1}{\pi(1+t)}\int_0^1 t^{-a}\sin(\pi a)\,da.
\]

For \(0<a<1\), \(t^{-a}\) is completely monotone because
\[
t^{-a}=\frac1{\Gamma(a)}\int_0^\infty e^{-ts}s^{a-1}\,ds.
\]
Also \(t\mapsto(1+t)^{-1}\) is completely monotone. Complete monotonicity is closed under products and positive mixtures, and \(\sin(\pi a)>0\) on \((0,1)\). Hence \(\phi_0\) is completely monotone.

If \(\phi\) is completely monotone and
\[
F(x)=\int_0^\infty e^{-xt}\phi(t)\,dt
\]
is finite for \(x>0\), then \(F\) is Stieltjes: write \(\phi(t)=\int_0^\infty e^{-ts}\,d\mu(s)\) and use Tonelli to obtain
\[
F(x)=\int_0^\infty\frac{d\mu(s)}{x+s}.
\]
Applying this to \(\phi_0\) proves that \(I_R\) is Stieltjes.

Finally, the source already proves that
\[
\widetilde I_R(x)=a+\int_0^\infty (1-e^{-xt})\phi_0(t)\,dt
\]
is a Bernstein function. The Levy density \(\phi_0\) is completely monotone, and it satisfies the Bernstein integrability condition
\[
\int_0^\infty (1\wedge t)\phi_0(t)\,dt<\infty.
\]
The standard complete-Bernstein criterion for Levy densities therefore gives that \(\widetilde I_R\) is a complete Bernstein function.
\end{proof}

The following appendix result records the fixed source problem \textbf{APP-0034} without interrupting the main exposition.
\begin{corollary}[APP-0034: Bulboaca-Zayed gamma-quotient monotonicity solved]
\label{res:bulboaca-zayed-gamma-quotient-full-monotonicity}
For the continuous extension \(\widetilde F\) of \(F(x)=\log\Gamma(x+1)/(\log(x^2+6)-\log(x+6))\) on \((-1,\infty)\), prove analytically that \(\widetilde F'\) has a unique zero \(x_m\simeq1.126207061\ldots\), with \(\widetilde F'<0\) on \((-1,x_m)\) and \(\widetilde F'>0\) on \((x_m,\infty)\).

This resolves the fixed application label \textbf{APP-0034}: Solve the Bulboaca-Zayed one-critical-point monotonicity statement for the gamma quotient.
\emph{Source reference.} \cite{app-app-0034-source}.
\end{corollary}
\begin{proof}
\emph{Setup.}
Use the notation from the previous Bulboaca--Zayed pass:
\[
G(x)=\log\Gamma(x+1),\qquad
D(x)=\log\frac{x^2+6}{x+6},
\]
and
\[
F(x)=\frac{G(x)}{D(x)}.
\]
For \(x>1\), \(D(x)>0\), and
\[
D'(x)=p(x)=\frac{x^2+12x-6}{(x^2+6)(x+6)}.
\]
Let
\[
N_0(x)=\psi(x+1)D(x)-G(x)D'(x).
\]
The derivative-sign numerator recorded in
the corresponding result is
\[
N(x)=((x^2+6)(x+6))\,N_0(x),
\]
so \(N\) and \(N_0\) have the same sign on \(x>1\).

The source itself proves the decreasing part on \((-1,1)\), and the earlier
the preceding theorem proves \(F'(x)>0\) for
\(x\ge8\).  It remains to prove a single sign change of \(N_0\) on \((1,8]\).

Put
\[
P(x)=x^2+12x-6,\qquad Q(x)=(x^2+6)(x+6),
\]
so \(p=P/Q\), and define
\[
r(x)=\frac{\psi(x+1)}{p(x)}=\psi(x+1)\frac{Q(x)}{P(x)}.
\]
Then
\[
P(x)^2r'(x)=S(x),
\]
where
\[
S(x)=\psi'(x+1)Q(x)P(x)+\psi(x+1)K(x),
\]
and
\[
K(x)=Q'(x)P(x)-Q(x)P'(x)
=x^4+24x^3+48x^2-144x-468.
\]
Direct differentiation gives
\[
S'(x)=P(x)\{Q(x)\psi''(x+1)+2Q'(x)\psi'(x+1)\}+K'(x)\psi(x+1).
\]

For \(y=x+1\ge2\), use the standard elementary bounds
\[
\psi(y)>\log y-\frac1y,\qquad
\psi'(y)>\frac1y+\frac{1}{2y^2},\qquad
\psi''(y)>-\frac1{y^2}-\frac2{y^3}.
\]
The last inequality follows from the series
\[
\psi''(y)=-2\sum_{n=0}^{\infty}(n+y)^{-3}
\]
and an integral upper bound for the sum.

Substituting these bounds into \(S'\) yields
\[
S'(x)>
P(x)\frac{5x^4+30x^3+57x^2+12x-90}{(x+1)^3}
K'(x)\left(\log(x+1)-\frac1{x+1}\right).
\]
On \(x\ge1\), \(P(x)>0\).  The quartic numerator has value \(14\) at \(x=1\)
and positive derivative
\[
20x^3+90x^2+114x+12.
\]
Also
\[
K'(x)=4x^3+72x^2+96x-144
\]
has value \(28\) at \(x=1\) and positive derivative for \(x\ge1\), while
\(\log(x+1)-1/(x+1)>0\) for \(x\ge1\).  Hence
\[
S'(x)>0\qquad(1\le x\le8).
\]

At the endpoints,
\[
S(1)=343\left(\frac{\pi^2}{6}-1\right)-539(1-\gamma)<0,
\]
for example using \(\pi^2<987/100\) and \(\gamma<5773/10000\), which gives
\(-66003/10000<0\).  Also \(S(8)>0\), because \(K(8)=17836>0\),
\(\psi'(9)>0\), and \(\psi(9)=H_8-\gamma>0\).  Therefore \(S\) has exactly
one zero on \([1,8]\).  Since \(P^2>0\), \(r'\) has exactly one zero; \(r\)
decreases first and then increases.

Now define
\[
I(x)=D(x)r(x)-G(x).
\]
Then
\[
N_0(x)=p(x)I(x),\qquad I'(x)=D(x)r'(x).
\]
Because \(D(x)>0\) for \(x>1\), \(I\) decreases and then increases with the
same turning point as \(r\).  Moreover \(\lim_{x\downarrow1}I(x)=0\), so
\(I(x)<0\) immediately to the right of \(1\).

Finally,
\[
N_0(8)=(H_8-\gamma)\log5-\frac{11}{70}\log(40320)>0.
\]
For instance, the bounds
\[
\gamma<\frac{5773}{10000},\quad \log5>\frac{1609}{1000},
\quad \log(40320)<\frac{10605}{1000}
\]
give the lower bound
\[
\left(\frac{761}{280}-\frac{5773}{10000}\right)\frac{1609}{1000}
-\frac{11}{70}\frac{10605}{1000}
=\frac{124435951}{70000000}>0.
\]
Thus \(I\), and therefore \(N_0\) and \(N\), has exactly one zero on
\((1,8]\).  Denote it by \(\xi\).  The sign pattern is
\[
N(x)<0\quad(1<x<\xi),\qquad
N(\xi)=0,\qquad
N(x)>0\quad(\xi<x\le8).
\]

A high-precision numerical check places
\[
\xi=1.12620706105164346558\ldots,
\]
matching the source's reported value.  The numerical value is not used in the
proof.

The raw numerator has a removable zero at \(x=1\), since \(G(1)=D(1)=0\).  Its
quadratic leading coefficient is
\[
\lim_{x\downarrow1}\frac{N_0(x)}{(x-1)^2}
=
\frac{7(\pi^2/6-1)-11(1-\gamma)}{98}<0,
\]
again by \(\pi^2<987/100\) and \(\gamma<5773/10000\), which gives
\(-1347/980000<0\).  Therefore any normalization dividing out the removable
\((x-1)^2\) factor is negative at the left endpoint.

The single-crossing candidate is true:
\[
T\text{-BZ-gamma-quotient-N-single-crossing-one-eight}.
\]
Together with the source's theorem on \((-1,1)\) and the previous local
right-tail theorem for \(x\ge8\), this proves the critical-window reduction
\[
T\text{-BZ-gamma-quotient-critical-window-reduction}
\]
and then the source problem
\[
T\text{-Bulboaca-Zayed-gamma-quotient-full-monotonicity}.
\]

The finite-envelope candidate is not separately promoted: the proof above is a
clean analytic ratio-kernel argument rather than the finite rational cover
specified in that candidate.  The diagnostic obstruction-map candidate is also
not needed after the single-crossing proof.
\end{proof}

The following appendix result records the fixed source problem \textbf{APP-0015} without interrupting the main exposition.
\begin{corollary}[APP-0015: Sibisi-Prabhakar measure is a stable-subordination push-forward]
\label{res:sibisi-prabhakar-q-measure-problem-solved-canonical}
In Sibisi's strict range \(0<\alpha<1\), \(\gamma>0\), \(\beta>\alpha\gamma\), the source problem of determining \(Q^\gamma_{\alpha,\beta}(\cdot\mid\lambda)\) with Laplace transform \(E^\gamma_{\alpha,\beta}(-\lambda x^\alpha)\) is solved at the canonical finite-measure level: \(Q\) is the stable subordination of the transform-normalized Pollard measure \(P^\gamma_{\alpha,\beta}\).

This resolves the fixed application label \textbf{APP-0015}: Give a canonical probabilistic representation of the Prabhakar Q-measure in the strict Pollard range.
\emph{Source reference.} \cite{app-app-0015-source}.
\end{corollary}
\begin{proof}
Conditioning on \(r\) and using the stable transform gives
\[
 \mathbb E e^{-x(\lambda r)^{1/\alpha}S_\alpha}=e^{-\lambda r x^\alpha}.
\]
Integration against \(P_{\alpha,\beta}^{\eta}\) gives the Pollard
representation of \(E_{\alpha,\beta}^{\eta}(-\lambda x^\alpha)\).  The total
mass is finite, equal to the value at \(x=0\), and no normalization to a
probability measure is needed for the transform identity.
\end{proof}
